\providecommand{\Kosz}{\mathsf{Kosz}}
\providecommand{\Conil}{\mathsf{Conil}}

\chapter{Theorem B: quadratic Koszul recognition}
\label{ch:theorem-B-scope-platonic}\label{ch:theorem-b-scope-platonic}

Let \(A=T(V)/(R)\) be a connected augmented quadratic chiral algebra.
Its quadratic coalgebra and universal bar coalgebra are distinct:
\[
 A^{\mathrm i}=C_X(s^{-1}V,s^{-2}R)\subset T_X^c(s^{-1}V),
 \qquad
 B_X(A)=T_X^c(s^{-1}\bar A).
\]
The quadratic twisting morphism
\(\tau_{\mathrm i}\colon A^{\mathrm i}\to A\) induces the canonical
comparison
\[
 q_A\colon A^{\mathrm i}\longrightarrow B_X(A).
\]
Theorem~B is the recognition theorem asserting that, under the chiral
Comparison Lemma package
\(H_{\mathrm{CL}}(A,A^{\mathrm i},\tau_{\mathrm i})\), the following
conditions are equivalent: \(q_A\) is a quasi-isomorphism;
\(\Omega_X(A^{\mathrm i})\to A\) is a quasi-isomorphism; the two
quadratic twisted tensor products are acyclic; and the PBW/bar
spectral sequence is strongly convergent and concentrated on the
quadratic Koszul diagonal.  Over a point this is Loday--Vallette
\cite[Theorems~2.3.2 and~3.4.6]{LV12}.  The chiral statement is
conditional precisely on \(H_{\mathrm{CL}}\).

Three neighbouring constructions supply the ambient in which this
comparison is studied.  The universal reconstruction counit
\[
 \varepsilon_A\colon\Omega_XB_X(A)\longrightarrow A
\]
belongs to Theorem~A and is independent of quadratic Koszulness
\cite[Corollary~2.3.4]{LV12}.  For a fixed chiral CDG-coalgebra \(C\),
Positselski's comodule--contramodule theorem compares
\(D^{\mathrm{co}}(C\text{-}\mathrm{CoFact})\) with
\(D^{\mathrm{ctr}}(C\text{-}\mathrm{ContraFact})\)
\cite[Theorem~5.2]{Positselski11}.  Finally, the category attached to
an acyclic twisting cochain is the separate construction
\(\mathsf{Tw}^{\mathrm{ch}}_{\mathrm{acyc}}(C,A,\tau)\) modeled on
Positselski~\cite[Theorem~6.5]{Positselski11}.  Thus the chapter
keeps four type signatures visible: universal reconstruction,
quadratic recognition, fixed-coalgebra second-kind categories, and
acyclic twisting modules.

Positselski's ordinary Correspondence Theorem~5.2 is intrinsic to an
arbitrary fixed CDG-coalgebra over a field.  Quadraticity and
Koszulness enter the separate twisting-morphism problem.  The
finite-window and completion packages below are sufficient hypotheses
for the chosen chiral de Rham transport and for assembling continuous
towers.

The organizing principle is the \emph{weight filtration}.  Assume the
chosen chiral algebra carries a positive-energy conformal grading
with finite-dimensional weight spaces.  Its bar coalgebra then
inherits a filtration by \emph{total bar weight}
$w = (\text{bar degree}) + \sum_i h_i$ summing the conformal weights
of the inputs. Each finite low-weight window
$F^{\leq w}\bar{B}^{\mathrm{ch}}(\cA)$ has finite-dimensional
state-space fibres in each cohomological degree.  A finite-dimensional
de Rham chain model additionally uses the holonomicity and de Rham
finiteness hypotheses of
Lemma~\ref{lem:finite-window-derham-realization-positselski}.  The
strict completed object is the finite-window tower together with its
continuous comparison data.  Its transition maps retain every
high-to-low contribution of the bar differential.

For the fixed-coalgebra comparison we isolate three statements.
First, at each finite weight, Positselski's ordinary fixed-coalgebra
comparison applies after de Rham realization; its chiral lift is the
additional package
\(\mathsf{Pos}^{\mathrm{ch}}_{\mathrm{co\text{-}ctr}}\)
(Theorem~\ref{thm:chiral-positselski-at-each-weight}). Second, in the
strict completed setting the continuous comparison package
\textup{(CP1)--(CP3)}
\textup{(}Definition~\ref{def:completed-chiral-positselski-package}\textup{)}
recovers the coderived/contraderived pair by the continuous adjoint functors
$\mathbb L\Phi^{\mathrm{ch}}_{\widehat C}$ and
$\mathbb R\Psi^{\mathrm{ch}}_{\widehat C}$
(Theorem~\ref{thm:chiral-positselski-weight-completed}); the unit,
counit, and contracting homotopies are inverse limits of finite-stage
data and live in the completed Hom complexes. Third, the direct sum
and its weight completion contain different state families
(Proposition~\ref{prop:chiral-positselski-raw-direct-sum-class-M-false}).
That linear-topological distinction determines which continuous
coactions and homotopies are defined in a chosen presentation.  The
universal reconstruction theorem remains unchanged.  A family-level
co--contra equivalence begins by constructing the fixed-coalgebra
hypotheses.

\section{Setup: the weight filtration on the chiral bar coalgebra}
\label{sec:weight-filtration-chiral-bar}

Let $X$ be a smooth complex curve, $\cA$ a chiral algebra on $X$ in
the sense of Beilinson--Drinfeld with conformal-weight grading
$\cA = \bigoplus_{h \geq 0} \cA_h$ and finite-dimensional weight
spaces $\dim \cA_h < \infty$. Write
$\bar{\cA} = \ker(\epsilon: \cA \to \mathbf{1})$ for the augmentation
ideal, and let
$\bar{B}^{\mathrm{ch}}(\cA) = T^c(s^{-1} \bar{\cA})$
denote the chiral bar coalgebra equipped with the deconcatenation
coproduct, bar differential inherited from the chiral product, and
the curved CDG structure
$(\bar{B}^{\mathrm{ch}}, d, h^{(g)})$. At genus~$0$ this is the
square-zero flat specialization used by the strict Koszul setting. At
positive genus the raw fibrewise identity is
\[
 d^{(g)\,2}=h^{(g)}\ast(-),
 \qquad
 m_1^{(g)\,2}(a)
 =
 m_2(m_0^{(g)},a)-m_2(a,m_0^{(g)})
 =
 [m_0^{(g)},a]_{m_2},
\]
in the transferred curved \(A_\infty\) notation.  The scalar
diagonal / uniform-weight projection sends the curvature class to
\[
 \operatorname{tr}_{\mathrm{diag}}\!\bigl(m_0^{(g)}\bigr)
 =
 \kappa(\cA)\,\omega_g.
\]
The period-corrected total modular differential is the square-zero
object, while the scalar display is its Hodge projection
(Theorem~\ref{thm:coalgebra-via-NAP}).

\begin{definition}[Chiral CDG comodule and contramodule categories;
\ClaimStatusDefinitional]
\label{def:chiral-cdg-co-contra-theorem-b}
Let $C$ be a chiral coalgebra on $X$ with differential $d_C$ and
curvature functional $h_C$. A \emph{chiral CDG-coalgebra} is the datum
$(C,d_C,h_C)$ with
\[
d_C^2=h_C\ast(-),\qquad h_C\circ d_C=0,
\]
where $\ast$ denotes the curvature action induced by the chiral
cooperations. A \emph{left chiral CDG-comodule} over $C$ is a chiral
comodule $(M,\rho_M)$ with differential $d_M$ compatible with $d_C$
and satisfying $d_M^2=h_C\ast(-)$. A \emph{left chiral
CDG-contramodule} over $C$ is a chiral contramodule $(P,\pi_P)$ with
differential $d_P$ compatible with $d_C$ and satisfying
$d_P^2(p)=h_C\cdot p$ through the contraaction.

Write $\operatorname{Hot}_{\mathrm{ch}}(C\text{-}\mathrm{comod})$
and
$\operatorname{Hot}_{\mathrm{ch}}(C\text{-}\mathrm{contra})$ for the
homotopy categories. Let
$\operatorname{Acycl}^{\mathrm{co}}_{\mathrm{ch}}(C)$ be the thick
triangulated subcategory generated by totalizations of short exact
sequences of chiral CDG-comodules and closed under the direct sums
admitted by the ambient category. Let
$\operatorname{Acycl}^{\mathrm{ctr}}_{\mathrm{ch}}(C)$ be the
analogous thick triangulated subcategory for chiral
CDG-contramodules, generated by short exact totalizations and closed
under the admitted products. The second-kind derived categories are
\[
D^{\mathrm{co}}_{\mathrm{ch}}(C)
:=
\operatorname{Hot}_{\mathrm{ch}}(C\text{-}\mathrm{comod})
/
\operatorname{Acycl}^{\mathrm{co}}_{\mathrm{ch}}(C),
\qquad
D^{\mathrm{ctr}}_{\mathrm{ch}}(C)
:=
\operatorname{Hot}_{\mathrm{ch}}(C\text{-}\mathrm{contra})
/
\operatorname{Acycl}^{\mathrm{ctr}}_{\mathrm{ch}}(C).
\]
The curvature identity \(d^2=h_C\ast(-)\) selects the second-kind
coacyclic and contraacyclic subcategories used in these localisations.
\end{definition}

\begin{definition}[Finite-stage chiral Positselski datum;
\ClaimStatusDefinitional]
\label{def:finite-stage-chiral-positselski-datum}
A finite stage $C_{\leq N}$ of a chiral CDG-coalgebra is a
\emph{finite-stage chiral Positselski datum} if:
\begin{enumerate}[label=\textup{(FP\arabic*)}]
\item each conformal-weight graded piece of $C_{\leq N}$ is
 finite-dimensional, and only finitely many weights occur in
 $C_{\leq N}$;
\item the reduced coproduct is conilpotent on $C_{\leq N}$, so the
 bar-length filtration terminates on every element;
\item the curvature $h_{\leq N}$ is a finite-stage chiral curvature:
 it is compatible with $d_{\leq N}$, satisfies
 $h_{\leq N}\circ d_{\leq N}=0$, and factors through the finite
 stage rather than through an infinite raw tail.
\end{enumerate}
Thus \textup{(FP1)} is the finite-dimensional graded-piece
condition, \textup{(FP2)} is conilpotence, and
\textup{(FP3)} is finite-stage curvature.
Theorem~\ref{thm:chiral-positselski-at-each-weight} additionally uses
the chiral lift
\(\mathsf{Pos}^{\mathrm{ch}}_{\mathrm{co\text{-}ctr}}\).
The completed theorem retains these hypotheses at every stage and
adds the tower package CP1--CP3 of
Definition~\ref{def:completed-chiral-positselski-package}.
\end{definition}

\begin{definition}[Total bar weight; \ClaimStatusDefinitional]
\label{def:total-bar-weight}
For a pure tensor
$\xi = s^{-1}\bar{a}_1 \otimes \cdots \otimes s^{-1}\bar{a}_n
 \in \bar{B}^{\mathrm{ch},n}(\cA)$
with $a_i \in \cA_{h_i}$, define the \emph{total bar weight} by
\begin{equation}\label{eq:total-bar-weight}
w(\xi) = n + \sum_{i=1}^{n} h_i.
\end{equation}
The \emph{weight filtration} $F^{\leq w} \bar{B}^{\mathrm{ch}}(\cA)$
is the subspace spanned by tensors of total bar weight at most $w$.
It is the finite low-weight window.  The conformal grading also splits
the underlying vector space into a complementary high-weight tail.
Since the bar differential may carry high-weight input to lower
weight, the finite chain object is
$F^{\leq w}\bar{B}^{\mathrm{ch}}(\cA)$ itself, or an explicitly chosen
finite-stage model carrying these high-to-low terms.
\end{definition}

\begin{lemma}[Filtration basics; \ClaimStatusProvedHere]
\label{lem:weight-filtration-basics}
The weight filtration satisfies:
\begin{enumerate}
\item $F^{\leq w}$ is preserved by the bar differential $d$ and by the
 deconcatenation coproduct $\Delta$;
\item each quotient
 $\mathrm{gr}^w = F^{\leq w}/F^{\leq w-1}$
 is finite-dimensional in each cohomological degree, with
 $\dim \mathrm{gr}^w_{n}
 = \sum_{n + \sum h_i = w} \prod_{i=1}^{n} \dim \cA_{h_i}$
 (a finite sum over the finite set of weak compositions of $w-n$ into
 nonnegative integers of bounded size);
\item the filtration is exhaustive:
 $\bigcup_{w \geq 0} F^{\leq w} = \bar{B}^{\mathrm{ch}}(\cA)$;
\item the associated graded
 $\mathrm{gr}^\bullet \bar{B}^{\mathrm{ch}}(\cA)
 = \bigoplus_{w \geq 0} \mathrm{gr}^w$
 is a conilpotent chiral coalgebra with finite-dimensional graded
 pieces; a filtration-preserving curvature equips it with the
 corresponding CDG structure.
\end{enumerate}
\end{lemma}

\begin{proof}
Items~(1) and~(3) are immediate from the combinatorial structure: the
bar differential and deconcatenation weakly decrease total bar weight,
and every tensor has some finite total weight by definition. Item~(2)
is finite-type because $\cA_h$ is finite-dimensional and only
finitely many partitions of $w$ into summands $(n; h_1,\ldots,h_n)$
with $\sum h_i = w - n$ contribute. Item~(4) follows from items~(1)
and~(2): on the associated graded, the iterated coproduct
$\bar{\Delta}^{(N)}$ vanishes on $\mathrm{gr}^{\leq N}$ because
deconcatenation strictly reduces bar degree when restricted to each
graded piece, and the finite-dimensionality in each weight provides
the conilpotency data.
\end{proof}

\begin{definition}[Strict weight-completed chiral bar object]
\label{def:weight-completed-bar}
The subcomplexes
\[
 F^{\leq0}\bar B^{\mathrm{ch}}(\cA)
 \hookrightarrow F^{\leq1}\bar B^{\mathrm{ch}}(\cA)
 \hookrightarrow\cdots
\]
form the exhaustive increasing weight filtration.  A
\emph{strict weight-completed chiral bar object} is the additional
inverse-system datum
\[
 \bigl(C_{\langle N\rangle},
 q_{N+1,N}\colon C_{\langle N+1\rangle}\twoheadrightarrow
 C_{\langle N\rangle}\bigr)_{N\geq0}
\]
of Lemma~\ref{lem:total-weight-rees-finite-cdg-tower}, together with
its comparison maps to the increasing filtration.  The differential,
curvature, coproduct, and chiral cooperations commute with the maps
\(q_{N+1,N}\), and
\[
 \widehat{\bar B}^{\mathrm{ch}}_{\mathrm{str}}(\cA)
 :=\varprojlim_N C_{\langle N\rangle}.
\]
For a family carrying algebra truncations
\(\cA_{\langle N\rangle}\) by chiral ideals, one may take
\(C_{\langle N\rangle}=\bar B^{\mathrm{ch}}
(\cA_{\langle N\rangle})\) once the induced transition maps satisfy
these conditions.
\end{definition}

\begin{remark}[Weight filtration on the state-space bar object]
\label{rem:why-weight-filtration}
Bar degree controls the number of tensor factors and witnesses
conilpotence on each finite word.  Total bar weight records the
conformal weights of genuine states.  For the universal Virasoro
vacuum module \(V_c\), put
\[
 v_k:=L_{-k}\mathbf1\in (V_c)_k,
 \qquad
 \mathbf e_k:=s^{-1}v_k\otimes s^{-1}v_k
 \in \bar B^{\mathrm{ch},2}(V_c),
 \qquad k\geq2.
\]
Then \(w(\mathbf e_k)=2k+2\).  A finite total-weight window contains
finitely many of the vectors \(\mathbf e_k\), and the weightwise
product completion contains their infinite sum.
\end{remark}

\begin{lemma}[A Virasoro state family in the weight completion;
\ClaimStatusProvedHere]
\label{lem:virasoro-mode-pair-no-raw-limit}
Let \(v_k=L_{-k}\mathbf1\) and
\(\mathbf e_k=s^{-1}v_k\otimes s^{-1}v_k\) for \(k\geq2\).  In the
underlying conformal-weight graded vector space set
\[
 V_\oplus:=\bigoplus_{k\geq2}\mathbb Q\mathbf e_k,
 \qquad
 V_\Pi:=\prod_{k\geq2}\mathbb Q\mathbf e_k.
\]
The family
\(
 \xi=(\mathbf e_2,\mathbf e_3,\ldots)
\)
defines an element of \(V_\Pi\) with infinite weight support.  The
image of \(V_\oplus\to V_\Pi\) is the finite-weight-support subspace,
so \(\xi\) lies in the complementary product subspace.
\end{lemma}

\begin{proof}
The vectors \(\mathbf e_k\) have distinct total weights \(2k+2\).
An element of \(V_\oplus\) has finitely many nonzero weight
components.  The displayed family has one nonzero component in every
weight \(2k+2\), and therefore belongs to \(V_\Pi\setminus V_\oplus\).
\end{proof}

\begin{proposition}[A completed comparison series and its discrete
support; \ClaimStatusProvedHere]
\label{prop:class-m-mode-family-obstruction-package}
Let \(E_\oplus=\bigoplus_{j\geq1}E_j\) and
\(E_\Pi=\prod_{j\geq1}E_j\).  Suppose a specified compatible family
of finite-window comparison maps or homotopies has value
\[
 h(x)=(e_1,e_2,\ldots)\in E_\Pi,
 \qquad e_j\in E_j,
\]
on some input \(x\), with infinitely many \(e_j\) nonzero.  This
particular comparison series belongs to the completed Hom complex,
and every representative of it in \(E_\oplus\) would have finite
support.  Hence its representing point lies in
\(E_\Pi\setminus\operatorname{im}(E_\oplus)\).
\end{proposition}

\begin{proof}
The canonical map \(E_\oplus\to E_\Pi\) has finite-support image.
The value \(h(x)\) has infinite support by hypothesis.
\end{proof}

\begin{remark}[The raw class-\(\mathsf M\) comparison problem;
\ClaimStatusOpen]
\label{rem:class-m-raw-comparison-open}
Let
\[
 \varepsilon_A^{\mathrm{raw}}\colon
 \Omega_X^{\mathrm{raw}}B_X^{\mathrm{raw}}(A)\longrightarrow A
\]
be a supplied discrete state-space realization of the chiral
bar--cobar counit.  The preceding proposition controls the chosen
finite-window comparison series: an infinite-support series belongs
to the completed Hom complex.  The existence of another discrete
homotopy is the contractibility problem for
\(\operatorname{Cone}(\varepsilon_A^{\mathrm{raw}})\) in the raw Hom
complex.  The class-\(\mathsf M\) cone remains uncomputed in this
chapter, and its ordinary quasi-isomorphism status is open.  The
completed comparison follows separately from
the explicit hypotheses \textup{(CP1)--(CP3)}.
\end{remark}

\section{The chiral Positselski equivalence at each finite weight}
\label{sec:chiral-positselski-at-each-weight}

\begin{theorem}[Finite-stage fixed-coalgebra comparison;
\ClaimStatusConditional]
\label{thm:chiral-positselski-at-each-weight}
\textup{[Type signature: Open quadrant; finite-window de Rham chiral
CDG-coalgebra presentation; Beilinson level~\(2\) module categories;
hypothesis package: a finite-stage CDG model together with
\(\mathsf{Pos}^{\mathrm{ch}}_{\mathrm{co\text{-}ctr}}\).]}
Let $\cA$ be a chiral algebra on a smooth curve $X$ with
positive-energy conformal-weight grading and finite-dimensional
weight spaces, and let
$C = \bar{B}^{\mathrm{ch}}(\cA)$ be its curved chiral CDG bar
coalgebra.  Assume the curvature action preserves total bar weight and
each \(F^{\leq w}C\) is a finite-stage chiral Positselski datum in the
sense of Definition~\ref{def:finite-stage-chiral-positselski-datum}.
Assume further that the finite-stage chiral lift
\(\mathsf{Pos}^{\mathrm{ch}}_{\mathrm{co\text{-}ctr}}
(F^{\leq w}C)\) of
Theorem~\ref{thm:chiral-bar-cobar-positselski-7-2} is supplied.
Then the ordinary de Rham realization carries Positselski's
fixed-coalgebra equivalence, and the supplied chiral lift gives
\begin{equation}\label{eq:chiral-positselski-at-each-weight}
D^{\mathrm{co}}\bigl((F^{\leq w}C)\text{-}
 \mathrm{comod}^{\mathrm{conil,\,ch}}\bigr)
\;\simeq\;
D^{\mathrm{ctr}}\bigl((F^{\leq w}C)\text{-}
\mathrm{contra}^{\mathrm{ch}}\bigr).
\end{equation}
Transition compatibility of these equivalences is an additional
assertion in \textup{(CP1)}.
\end{theorem}

\begin{proof}
Lemma~\ref{lem:weight-filtration-basics}(1),(2),(4), together with the
curvature hypothesis, shows that \(F^{\leq w}C\) is a
differential-stable chiral CDG-subcoalgebra and is finite-dimensional
in each cohomological degree
(a finite direct sum of finite-dimensional graded pieces). The
conilpotency is witnessed by the bar-degree-plus-weight filtration:
on any element of $F^{\leq w}C$, the iterated deconcatenation
coproduct $\bar{\Delta}^{(N)}$ vanishes once $N$ exceeds the maximum
bar degree in $F^{\leq w}C$, which is bounded by~$w$ itself (since
$n \leq w$ when $n + \sum h_i \leq w$ with $h_i \geq 0$).

After de Rham realization, Positselski's Theorem~5.2 applies to the
resulting ordinary CDG-coalgebra over the ground field.  That theorem
allows arbitrary CDG-coalgebras; finite dimensionality and
conilpotence here serve the finite-window construction rather than the
ordinary fixed-coalgebra theorem.  The package
\(\mathsf{Pos}^{\mathrm{ch}}_{\mathrm{co\text{-}ctr}}\) lifts the two
functors, their unit and counit, and the second-kind acyclic
subcategories back to the chiral finite-stage model.  This gives
\eqref{eq:chiral-positselski-at-each-weight}.
\end{proof}

\begin{remark}[Family input at finite weight]
\label{rem:independence-from-boundary}
For every family, a finite-window application consists of a
differential-stable CDG model and the chiral lift
\(\mathsf{Pos}^{\mathrm{ch}}_{\mathrm{co\text{-}ctr}}\).
Positive-energy finite dimensionality supplies the underlying
finite-weight vector spaces.  The chiral functors, preservation of
second-kind acyclics, and factorization compatibility remain explicit
family inputs.
\end{remark}

\section{The weight-completed chiral Positselski equivalence}
\label{sec:chiral-positselski-weight-completed}

\begin{definition}[Strict Mittag--Leffler completed chiral coalgebra;
\ClaimStatusDefinitional]
\label{def:strict-ml-completed-chiral-coalgebra}
A \emph{strict Mittag--Leffler completed chiral CDG-coalgebra} is an
inverse system $\{C_N\}_{N\geq0}$ of finite-window de Rham chiral
CDG-coalgebras such that:
\begin{enumerate}[label=\textup{(ML\arabic*)}]
\item every $C_N$ is conilpotent and finite-dimensional in each
cohomological degree;
\item the transition maps $C_{N+1}\twoheadrightarrow C_N$ are strict
surjective morphisms of chiral CDG-coalgebras;
\item the differential, curvature, coproduct, and chiral cooperations
commute with the transition maps and define continuous operations on
$\widehat C:=\varprojlim_N C_N$;
\item for every cohomological degree $m$, the inverse systems
$\{H^m(C_N)\}_N$ and the corresponding cone systems are
Mittag--Leffler.
\end{enumerate}
Continuous comodules and contramodules over $\widehat C$ are compatible
inverse systems of finite-stage comodules and contramodules. Their
morphism complexes are the derived inverse limits of finite-stage
morphism complexes.
\end{definition}

\begin{definition}[Completed chiral Positselski package CP1--CP3;
\ClaimStatusDefinitional]
\label{def:completed-chiral-positselski-package}
Let
\[
\widehat C=\varprojlim_N C_N
\]
be a strict Mittag--Leffler completed chiral CDG-coalgebra in the sense
of Definition~\ref{def:strict-ml-completed-chiral-coalgebra}. The
\emph{completed chiral Positselski package}
\(\mathsf{CP}^{\mathrm{ch}}(\widehat C)\) consists of the following
three assertions.
\begin{enumerate}
\item[\textup{(CP1)}] The finite stages \(C_N\) are finite-stage
 chiral Positselski data in the sense of
 Definition~\ref{def:finite-stage-chiral-positselski-datum}: they have
 finite-dimensional graded pieces, terminating conilpotence, and
 finite-stage curvature.  Each stage carries
 \(\mathsf{Pos}^{\mathrm{ch}}_{\mathrm{co\text{-}ctr}}(C_N)\), and
 the transition functors intertwine the finite-stage co--contra
 adjunctions, their units, and their counits.  The transition maps
 \(C_{N+1}\twoheadrightarrow C_N\) are strict surjective chiral
 CDG-coalgebra morphisms, and the differential, curvature, coproduct,
 and chiral cooperations define continuous operations on
 \(\widehat C\).
\item[\textup{(CP2)}] The completed coderived category is the strict
 derived inverse limit of the finite-stage coderived categories:
 \[
 D^{\mathrm{co}}_{\mathrm{ch}}(\widehat C)
 \simeq
 R\!\varprojlim_N D^{\mathrm{co}}_{\mathrm{ch}}(C_N).
 \]
 Equivalently, Verdier localization by the coacyclic subcategories may
 be performed before or after the strict derived inverse limit. The
 Hom-complex and cone towers that occur in units, counits, and
 comparison maps are Mittag--Leffler, so the relevant
 \(\varprojlim^1\) terms vanish.
\item[\textup{(CP3)}] The completed contraderived category is similarly
 the strict derived inverse limit:
 \[
 D^{\mathrm{ctr}}_{\mathrm{ch}}(\widehat C)
 \simeq
 R\!\varprojlim_N D^{\mathrm{ctr}}_{\mathrm{ch}}(C_N).
 \]
 The completed chiral co/contra adjunction preserves the
 product-generated contraacyclic objects and the sum-generated
 coacyclic objects. A transition-compatible finite-stage family of
 units, counits, or contracting homotopies therefore has a continuous
 inverse limit in the completed Hom complex.  In a split weight
 product, this particular inverse-limit element belongs to the image
 of the discrete Hom complex exactly when it has finite weight
 support.
\end{enumerate}
This package is a statement about the strict inverse-limit ambient and
retains the transition maps and topology as part of \(\widehat C\).
\end{definition}

\begin{lemma}[Strict total-weight Rees tower criterion;
\ClaimStatusProvedHere]
\label{lem:total-weight-rees-finite-cdg-tower}
Let $C=\bar B^{\mathrm{ch}}(\cA)$ carry the total-weight filtration
of Definition~\ref{def:total-bar-weight}.  A strict total-weight
Rees tower for $C$ is the following additional datum:
\[
\bigl(C_{\langle N\rangle},q_{N+1,N},j_N,r_N\bigr)_{N\geq0}.
\]
Here $C_{\langle N\rangle}$ is a finite-window chiral
CDG-coalgebra, $q_{N+1,N}\colon C_{\langle N+1\rangle}\twoheadrightarrow
C_{\langle N\rangle}$ is a strict surjective CDG-coalgebra morphism,
and
\[
j_N\colon F^{\leq N}C\longrightarrow C_{\langle N\rangle},
\qquad
r_N\colon C_{\langle N\rangle}\longrightarrow F^{\leq N}C
\]
are filtered comparison maps inducing inverse isomorphisms on
$\operatorname{gr}^w_F$ for every $w\leq N$. The compatibility
condition is
\[
q_{N+1,N}\circ j_{N+1}\circ \iota_{N,N+1}
=
j_N
\quad\text{and}\quad
r_N\circ q_{N+1,N}
=
\pi^{\mathrm{low}}_{N+1,N}\circ r_{N+1}
\]
on the low-weight output window, where
$\iota_{N,N+1}\colon F^{\leq N}C\hookrightarrow F^{\leq N+1}C$ is
inclusion and $\pi^{\mathrm{low}}_{N+1,N}$ is the comparison map
specified by the finite-window model. Then
\[
\widehat C_{\mathrm{str}}
:=\varprojlim_N C_{\langle N\rangle}
\]
is a strict completed chiral CDG-coalgebra whose operations recover
the total-weight-completed operations on $C$ through the maps
$j_N,r_N$.

The datum of the maps $q_{N+1,N}$ is essential.  The projection
\[
\bigoplus_{0\leq w\leq N+1}\operatorname{gr}^w_F C
\longrightarrow
\bigoplus_{0\leq w\leq N}\operatorname{gr}^w_F C
\]
is a chain map precisely when the chiral bar differential has zero
component from weight \(N+1\) to lower weights.  In general the strict
tower uses typed transition maps that retain this component.
\end{lemma}

\begin{proof}
Lemma~\ref{lem:weight-filtration-basics} gives the filtered
finite-type pieces.  The hypotheses on
$C_{\langle N\rangle},q_{N+1,N},j_N,r_N$ assert exactly the missing
chain-level datum: a finite model whose transition maps commute with
the full filtered differential, curvature, coproduct, and chiral
cooperations.  Finite-dimensionality follows from the associated
graded comparison with $\operatorname{gr}^{\leq N}_F C$.  Strict
surjectivity of the $q_{N+1,N}$ gives a genuine inverse system of
CDG-coalgebras, and the inverse limit carries the unique continuous
operations detected on all finite stages.  The final paragraph is the
direct calculation showing that the typed transition maps carry every
weight-lowering component of the differential.
\end{proof}

\begin{lemma}[Finite-window de Rham realisation for Positselski;
\ClaimStatusConditional]
\label{lem:finite-window-derham-realization-positselski}
At every strict total-weight Rees stage $C_{\langle N\rangle}$ of
Lemma~\ref{lem:total-weight-rees-finite-cdg-tower}, de Rham global
sections on the fixed curve realise the chiral CDG-coalgebra,
comodules, and contramodules in $\mathrm{Ch}(\mathrm{Vect})$ with
finite-dimensional graded pieces in each cohomological degree,
provided the finite-stage \(\mathcal D_X\)-modules are holonomic and
their de Rham cohomology is finite-dimensional. In
this finite-window de Rham ambient, the exactness hypotheses in
Positselski's second-kind comodule--contramodule correspondence hold:
totalisations of short exact triples generate coacyclic objects,
dual product-totalisations generate contraacyclic objects, and finite
direct sums and products are exact.
\end{lemma}

\begin{proof}
The strict Rees stage has finite-dimensional conformal-weight fibres
by its associated-graded comparison with
$\operatorname{gr}^{\leq N}_F C$ and
Lemma~\ref{lem:weight-filtration-basics}(2).  The holonomicity and
de Rham finiteness hypotheses give finite-dimensional global
cohomology in each fixed degree and window.  After de Rham global
sections, chiral operations become ordinary operations on complexes
of vector spaces. In $\mathrm{Ch}(\mathrm{Vect})$ kernels,
cokernels, sums, and products are computed degreewise, and vector
spaces over a field are exact for both sums and products. Thus the
finite-stage totalisations used to define the second-kind
coacyclic and contraacyclic subcategories satisfy Positselski's
exactness hypotheses \cite[\S\S3--4]{Positselski11}.
\end{proof}

\begin{lemma}[Hom-cone Mittag--Leffler for strict finite-window towers;
\ClaimStatusProvedHere]
\label{lem:hom-cone-ml-rees-towers}
Let $\{E_N\}_N$ and $\{F_N\}_N$ be compatible towers of finite-stage
complexes arising from the strict finite-window tower of
Lemma~\ref{lem:total-weight-rees-finite-cdg-tower}, and let
$f_N\colon E_N\to F_N$ be a compatible system of chain maps.  Suppose
that for every cohomological degree $m$ the inverse systems
$\{H^m(E_N)\}_N$, $\{H^m(F_N)\}_N$, and
$\{H^m\operatorname{Hom}(E_N,F_N)\}_N$ are Mittag--Leffler. Then
$\{H^m(\operatorname{Cone}(f_N))\}_N$ is Mittag--Leffler for every
$m$, and the Hom tower has the Milnor identification
\[
H^m R\!\varprojlim_N\operatorname{Hom}(E_N,F_N)
\cong
\varprojlim_N H^m\operatorname{Hom}(E_N,F_N).
\]
Moreover, the cone tower has the Milnor exact sequence
\[
0\to
\varprojlim\nolimits^1_N
H^{m-1}\!\bigl(\operatorname{Cone}(f_N)\bigr)
\to
H^m\!\bigl(R\!\varprojlim_N\operatorname{Cone}(f_N)\bigr)
\to
\varprojlim_N
H^m\!\bigl(\operatorname{Cone}(f_N)\bigr)
\to0 .
\]
In particular, if each $f_N$ is a quasi-isomorphism, then
$R\!\varprojlim_N f_N$ is a quasi-isomorphism.
\end{lemma}

\begin{proof}
For each $N$ the cone fits into the standard distinguished triangle
$E_N\to F_N\to\operatorname{Cone}(f_N)\to E_N[1]$.  Taking cohomology
gives a compatible long exact sequence of inverse systems.  In the
strict finite-window tower all cohomology groups are finite-dimensional
vector spaces at each fixed stage and degree.  Images in an inverse
system of finite-dimensional vector spaces stabilize exactly when the
Mittag--Leffler condition holds; the long exact sequence then shows
that the cone cohomology systems are Mittag--Leffler. The Hom
identification is the Milnor sequence for the tower of Hom complexes,
with $\varprojlim^1$ killed by the stated ML hypothesis and
\cite[Prop.~3.5.7]{Weibel94}. The displayed cone sequence is the
same Milnor sequence applied to the tower of cones. If all $f_N$ are
quasi-isomorphisms, the cone cohomology systems vanish, so both the
inverse-limit and $\varprojlim^1$ terms vanish, and the derived
inverse-limit cone is acyclic.
\end{proof}

\begin{proposition}[Verdier localization through a strict inverse
limit; \ClaimStatusConditional]
\label{prop:strict-rlim-verdier-quotient}
\label{lem:strict-rlim-verdier-quotient}
Let \(\{\mathcal T_N\}_N\) be a strict inverse tower of finite-window
dg categories and
\(\mathcal A_N\subset\mathcal T_N\) compatible thick triangulated dg
subcategories.  Assume the localization compatibility clause
\textup{(CP2)}: localized mapping complexes satisfy the Hom-cone
Mittag--Leffler condition, every strict compatible target tower has a
finite-stage representative, and the canonical functor is fully
faithful on those representatives.  Then
\[
 R\!\varprojlim_N(\mathcal T_N/\mathcal A_N)
 \xrightarrow{\sim}
 \bigl(R\!\varprojlim_N\mathcal T_N\bigr)/
 \bigl(R\!\varprojlim_N\mathcal A_N\bigr)
\]
on the full subcategory of strict compatible towers.
\end{proposition}

\begin{proof}
Full faithfulness and representability are the two clauses of
\textup{(CP2)}.  They give full faithfulness and essential
surjectivity of the displayed functor.
\end{proof}

\begin{proposition}[Second-kind acyclicity under strict completion;
\ClaimStatusConditional]
\label{prop:completed-second-kind-generator-preservation}
Let \(\{C_N\}_N\) be a strict finite-window tower.  Assume the
second-kind compatibility clause \textup{(CP3)}: strict inverse limit
commutes with the admitted direct sums, products, and short-exact
totalizations, and the finite-stage co--contra functors preserve their
coacyclic and contraacyclic generators.  Then the completed functors
preserve the completed second-kind acyclic subcategories and descend
to the coderived and contraderived categories.
\end{proposition}

\begin{proof}
The generators are totalizations of short exact sequences, followed
by the admitted sums on the comodule side and products on the
contramodule side.  Clause \textup{(CP3)} carries these generators
through the strict inverse limit and through the completed adjunction.
Verdier localization therefore produces the asserted functors.
\end{proof}

\begin{theorem}[Continuous fixed-coalgebra co--contra comparison;
\ClaimStatusConditional]
\label{thm:chiral-positselski-weight-completed}
\phantomsection\label{thm:chiral-positselski-7-2}%
\phantomsection\label{thm:chiral-positselski-5-3}%
\textup{[Type signature: Open quadrant; strict inverse-limit
fixed-coalgebra presentation; Beilinson level~\(2\) module
categories; hypothesis package:
\(\mathsf{CP}^{\mathrm{ch}}(\widehat C)\).]}
Let
\(\widehat C=\varprojlim_NC_N\) be a strict Mittag--Leffler completed
chiral CDG-coalgebra.  Assume the completed chiral Positselski package
\(\mathsf{CP}^{\mathrm{ch}}(\widehat C)\):
\textup{(CP1)} gives finite-stage fixed-coalgebra comparisons and
continuous structure maps; \textup{(CP2)} identifies the completed
coderived category with the derived inverse limit; and
\textup{(CP3)} gives the corresponding contraderived identification
and preservation of second-kind acyclic objects.  Then
\begin{equation}\label{eq:chiral-positselski-weight-completed}
D^{\mathrm{co}}\bigl(\widehat C\text{-}
 \mathrm{comod}^{\mathrm{conil,ch}}\bigr)
\simeq
R\!\varprojlim_N
D^{\mathrm{co}}\bigl(C_N\text{-}
 \mathrm{comod}^{\mathrm{conil,ch}}\bigr)
\simeq
R\!\varprojlim_N
D^{\mathrm{ctr}}\bigl(C_N\text{-}
 \mathrm{contra}^{\mathrm{ch}}\bigr)
\simeq
D^{\mathrm{ctr}}\bigl(\widehat C\text{-}
 \mathrm{contra}^{\mathrm{ch}}\bigr).
\end{equation}
The equivalence is realized by the continuous adjoint functors
\[
 \widehat\Phi^{\mathrm{ch}}
 =R\!\varprojlim_N\Phi_N^{\mathrm{ch}},
 \qquad
 \widehat\Psi^{\mathrm{ch}}
 =R\!\varprojlim_N\Psi_N^{\mathrm{ch}}.
\]
Compatible finite-stage units, counits, and contracting homotopies
define their continuous inverse-limit counterparts.  Such a homotopy
belongs to the completed Hom complex.  In a split weight product, its
membership in the discrete Hom complex is equivalent to finite weight
support.
\end{theorem}

\begin{proof}
The finite-stage fixed-coalgebra equivalences form a natural tower by
\textup{(CP1)}.  Taking \(R\!\varprojlim\) gives the middle
equivalence.  The identifications in \textup{(CP2)} and
\textup{(CP3)} supply the outer equivalences and ensure that the
finite-stage coacyclic and contraacyclic cones remain acyclic after
completion.  The Hom-cone Mittag--Leffler condition identifies the
cone of an inverse-limit unit or counit with the inverse limit of its
finite-stage cones.  The same argument sends a compatible family of
finite-stage homotopies to a continuous homotopy.
\end{proof}

\begin{remark}[Relation to the MC4 completion theorem]
\label{rem:relation-to-MC4-completion}
The weight-completed equivalence
\eqref{eq:chiral-positselski-weight-completed} and
Theorem~\ref{thm:completed-bar-cobar-strong} use the same inverse-limit
technology in different type signatures.  MC4 is a chain-level map in
the associative-algebra category.  The present theorem is an
equivalence between second-kind module categories over one fixed
coalgebra.  Each proof has its own cone tower and Mittag--Leffler
obligation.
\end{remark}

\section{Direct-sum and completed state families}
\label{sec:raw-direct-sum-class-M-false}

\begin{proposition}[State-family separation for the class-\(\mathsf M\)
presentation; \ClaimStatusProvedHere]
\label{prop:chiral-positselski-raw-direct-sum-class-M-false}
For the universal Virasoro vacuum module \(V_c\), let
\[
 v_k=L_{-k}\mathbf1,\qquad
 \mathbf e_k=s^{-1}v_k\otimes s^{-1}v_k,\qquad k\geq2 .
\]
The total weight of \(\mathbf e_k\) is \(2k+2\).  The direct sum and
its split weight completion contain respectively
\[
 V_\oplus=\bigoplus_{k\geq2}\mathbb Q\mathbf e_k,
 \qquad
 V_\Pi=\prod_{k\geq2}\mathbb Q\mathbf e_k .
\]
The family \(\xi=(\mathbf e_2,\mathbf e_3,\ldots)\) lies in
\(V_\Pi\setminus V_\oplus\).  Consequently, a specified compatible
finite-window homotopy whose value contains these components belongs
to the completed Hom complex and has infinite weight support.  This
conclusion concerns that compatible homotopy.  The ordinary-derived
status of the raw counit is the open cone problem in
Remark~\ref{rem:class-m-raw-comparison-open}.
\end{proposition}

\begin{proof}
Apply Lemma~\ref{lem:virasoro-mode-pair-no-raw-limit} and
Proposition~\ref{prop:class-m-mode-family-obstruction-package}.
\end{proof}

\begin{remark}[Universal reconstruction for the standard families;
 \ClaimStatusProvedElsewhere]
\label{rem:theorem-B-chain-level-G-L-attribution}
The ordinary universal counit
\(\Omega B(A)\to A\) is the quasi-free bar resolution of every
connected augmented dg algebra
\cite[Corollary~2.3.4]{LV12}.  Its chiral realization requires the
Ran comparison and factorization package of Theorem~A.  Family names
enter after construction of that transport package, while
quadratic Koszulness enters through the separate map
\(q_A\colon A^{\mathrm i}\to B_X(A)\).

The label \(\texttt{thm:bv-bar-coderived-vol1}\) belongs to Vol.~I,
Theorem~\ref{thm:bv-bar-coderived-vol1}, in
\texttt{chapters/connections/bv\_brst.tex}.  Its status is
\emph{conditional}, with the harmonic-factorization and
coderived hypotheses stated there.  Vol.~II cites this Vol.~I result.
That BV/bar comparison and the completed fixed-coalgebra assertions
\textup{(CP1)--(CP3)} have distinct source and target categories.
\end{remark}

\begin{remark}[Bar-length conilpotence and weight completion]
\label{rem:class-G-L-survive-raw}
The cofree bar coalgebra is conilpotent by bar length on every finite
bar word.  A weight completion controls infinite compatible state
families and continuous operations.  A fixed-coalgebra co--contra
comparison therefore records both properties separately: algebraic
conilpotence of the coalgebra and continuity/exactness in the chosen
topological presentation.
\end{remark}

\begin{definition}[Feigin--Frenkel centre weight completion;
\ClaimStatusDefinitional]
\label{def:ff-center-weight-completion-theorem-b}
\index{Feigin--Frenkel centre!weight completion}
Let \(\fg\) be simple of rank \(r\), with exponents
\(d_1,\ldots,d_r\), and let
\(\mathfrak z(\widehat{\fg})\) denote the Feigin--Frenkel centre of
the critical affine algebra. On the formal disc, write its
Segal--Sugawara generators as
\[
S_i^{(m)}:=\partial^m S_i,\qquad
\operatorname{wt}(S_i^{(m)})=d_i+1+m,\qquad
1\leq i\leq r,\; m\geq0 .
\]
The \emph{Feigin--Frenkel weight filtration} is
\[
F^{\leq N}\mathfrak z(\widehat{\fg})
:=
\bigoplus_{0\leq w\leq N}\mathfrak z(\widehat{\fg})_w,
\qquad
F^{>N}\mathfrak z(\widehat{\fg})
:=
\bigoplus_{w>N}\mathfrak z(\widehat{\fg})_w .
\]
The finite-window centre and its weight completion are
\[
\mathfrak z_{\leq N}(\widehat{\fg})
:=
\mathfrak z(\widehat{\fg})/
F^{>N}\mathfrak z(\widehat{\fg}),
\qquad
\widehat{\mathfrak z}_{\mathrm{wt}}(\widehat{\fg})
:=
\varprojlim_N \mathfrak z_{\leq N}(\widehat{\fg})
\cong
\prod_{w\geq0}\mathfrak z(\widehat{\fg})_w .
\]
The neighbourhood basis of zero is the tail system
\(\prod_{w>N}\mathfrak z(\widehat{\fg})_w\). This is the completion
topology used for the Feigin--Frenkel-centre clause of
Corollary~\ref{cor:positselski-applicable-families}.  It supplies the
topological part of \textup{(CP1)} for the commutative centre tower;
the coderived and contraderived assertions \textup{(CP2)}--\textup{(CP3)}
remain separate inputs in
Definition~\ref{def:completed-chiral-positselski-package}.
\end{definition}

\begin{lemma}[Finite Feigin--Frenkel weight windows;
\ClaimStatusProvedHere]
\label{lem:ff-center-finite-weight-windows}
\index{Feigin--Frenkel centre!finite windows}
For the centre with the weight filtration of
Definition~\ref{def:ff-center-weight-completion-theorem-b}, every
weight piece \(\mathfrak z(\widehat{\fg})_w\) is finite-dimensional,
every finite window
\(\mathfrak z_{\leq N}(\widehat{\fg})\) is finite-dimensional, and the
transition maps
\(\mathfrak z_{\leq N+1}(\widehat{\fg})\twoheadrightarrow
\mathfrak z_{\leq N}(\widehat{\fg})\) are strict surjections. The
commutative product and translation operator \(\partial\) are continuous
for this topology. These facts establish the finite-window and continuity
part of \textup{(CP1)} for the centre tower.  Second-kind acyclicity
and Verdier-localisation are the additional assertions
\textup{(CP2)}--\textup{(CP3)}.
\end{lemma}

\begin{proof}
By the Feigin--Frenkel theorem, on the formal disc
\(\mathfrak z(\widehat{\fg})\) is the polynomial algebra on the
generators \(S_i^{(m)}\) above; this is the same generator statement
recorded in Theorem~\ref{thm:critical-level-structure}. Fix \(w\). A
monomial of total conformal weight \(w\) can involve only generators with
\(d_i+1+m\leq w\). There are finitely many such pairs \((i,m)\), and a
finite set of positive generator weights has only finitely many
monomials of total weight \(w\). Hence
\(\dim \mathfrak z(\widehat{\fg})_w<\infty\), and summing over
\(0\leq w\leq N\) gives finite-dimensionality of
\(\mathfrak z_{\leq N}\).

Weight is additive under multiplication, so \(F^{>N}\mathfrak z\) is an
ideal and the quotient maps form a strict surjective inverse system.
Multiplication is continuous because the product of tails lands in the
corresponding additive tail, and \(\partial\) is continuous because
\(\partial S_i^{(m)}=S_i^{(m+1)}\) raises weight by one.  This proves
the stated finite-window topology and continuity assertions.
\end{proof}

\section{Family hypotheses for the completed comparison}
\label{sec:family-by-family-applicability}

\begin{corollary}[Hypotheses for a fixed-coalgebra
co--contra comparison; \ClaimStatusConditional]
\label{cor:positselski-applicable-families}
For each of the Heisenberg, affine Kac--Moody, free-field, Virasoro,
\(W_N\), Yangian, and Feigin--Frenkel families, applicability of the
fixed-coalgebra chiral Positselski comparison requires the following
data:
\begin{enumerate}[label=\textup{(F\arabic*)}]
\item a specified chiral CDG-coalgebra \(C\);
\item a de Rham realization and the chiral lift
 \(\mathsf{Pos}^{\mathrm{ch}}_{\mathrm{co\text{-}ctr}}(C)\),
 including preservation of the second-kind acyclic subcategories;
\item for an inverse-limit presentation, the completed package
 \textup{(CP1)--(CP3)} of
 Definition~\ref{def:completed-chiral-positselski-package}.
\end{enumerate}
These are applicability criteria rather than family-name
consequences.  For the Feigin--Frenkel centre,
Definition~\ref{def:ff-center-weight-completion-theorem-b} and
Lemma~\ref{lem:ff-center-finite-weight-windows} establish the topology
and finite-window dimensions.  The completed co--contra comparison
applies once this tower also satisfies \textup{(CP2)} and
\textup{(CP3)}.

The worked calculations in this chapter establish the finite
weight dimensions and the topological part of \textup{(CP1)} for the
Feigin--Frenkel centre.  They leave \textup{(CP2)} and \textup{(CP3)}
as open categorical construction problems.  The other named families
likewise retain their stated family inputs.  Thus every family-level
completed chiral co--contra comparison in this chapter has
\emph{conditional} status.
\end{corollary}

\begin{proof}
Positselski's theorem is functorial in the fixed coalgebra and uses
the stated second-kind exactness hypotheses.  The completed statement
is obtained from the finite-window theorem after the derived
inverse-limit and localization comparisons in \textup{(CP1)--(CP3)}.
Conilpotence enters the separate Loday--Vallette bar--cobar unit and
counit and the quadratic recognition map \(q_A\).  Positselski's
ordinary fixed-coalgebra correspondence applies to every CDG
coalgebra over the ground field.
\end{proof}

\begin{remark}[Separation of the four scopes]
\label{rem:unconditional-scope-clarification}
The direct-sum/completion distinction describes a topological
presentation.  Fixed-coalgebra co--contra equivalence is a
module-category theorem.  Universal bar--cobar reconstruction is
Theorem~A.  Quadratic recognition is Theorem~B and is detected by
\(q_A\colon A^{\mathrm i}\to B_X(A)\).
\end{remark}

\section{Four typed comparisons at Beilinson levels
\(1\leftrightarrow2\)}
\label{sec:scope-separation-thm-B}

The chapter contains four constructions with four distinct type
signatures:
\[
\begin{array}{c|c|c}
\text{comparison} & \text{canonical datum} & \text{mathematical source}\\ \hline
\text{universal reconstruction} &
 \varepsilon_A\colon\Omega_XB_X(A)\to A &
 \text{Theorem~A; \cite[Corollary~2.3.4]{LV12}}\\
\text{quadratic recognition} &
 q_A\colon A^{\mathrm i}\to B_X(A) &
 \text{Theorem~B; \cite[Theorem~3.4.6]{LV12}}\\
\text{fixed-coalgebra co--contra} &
 D^{\mathrm{co}}(C\text{-}\mathrm{CoFact})
 \simeq D^{\mathrm{ctr}}(C\text{-}\mathrm{ContraFact}) &
 \text{\cite[Theorem~5.2]{Positselski11}}\\
\text{acyclic twisting modules} &
 \mathsf{Tw}^{\mathrm{ch}}_{\mathrm{acyc}}(C,A,\tau) &
 \text{\cite[Theorem~6.5]{Positselski11}} .
\end{array}
\]
The first two rows compare algebra and coalgebra objects.  The third
compares module categories over one fixed coalgebra.  The fourth
compares module categories attached to a supplied acyclic twisting
cochain.  This typing determines the hypotheses and the meaning of
every equivalence in the chapter.

\begin{remark}[The sharp Theorem~B statement;
\ClaimStatusConditional]
\label{rem:thm-B-scope-confusion-type}
\textup{[Type signature: Open quadrant; connected complete quadratic
chiral presentation; Beilinson levels \(1\leftrightarrow2\);
hypothesis package
\(H_{\mathrm{CL}}(A,A^{\mathrm i},\tau_{\mathrm i})\).]}
Theorem~B concerns
\[
 A^{\mathrm i}=C_X(s^{-1}V,s^{-2}R)
 \xrightarrow{\,q_A\,} B_X(A)
\]
for \(A=T_X(V)/(R)\).  Its status is the status of \(q_A\).
Completion enters through convergence of this comparison.  Verdier
duality enters through the transpose
\(\mathbb D_{\operatorname{Ran}}(q_A)\).  Fixed-coalgebra second-kind
categories provide a neighbouring categorical construction with
their own hypothesis package.
\end{remark}

\section{The bar--cobar adjunction and universal reconstruction}
\label{sec:thm-B-coalgebra-side-unit-qi}

\begin{lemma}[Canonical direction of the chiral bar--cobar adjunction;
\ClaimStatusConditional]
\label{lem:bar-cobar-adjunction-direction-canonical}
For a conilpotent chiral coalgebra \(C\) and an augmented chiral
algebra \(A\), the chiral twisting-morphism dictionary gives
\begin{equation}\label{eq:bar-cobar-hom-bijection-canonical}
\operatorname{Hom}_{\mathrm{Alg}^{\mathrm{ch,aug}}_X}
 \bigl(\Omega_XC,A\bigr)
\cong
\operatorname{Tw}^{\mathrm{ch}}(C,A)
\cong
\operatorname{Hom}_{\mathrm{CoAlg}^{\mathrm{ch,conil}}_X}
 \bigl(C,B_X(A)\bigr).
\end{equation}
Hence \(\Omega_X\dashv B_X\), with unit
\(\eta_C\colon C\to B_X\Omega_XC\) and counit
\(\varepsilon_A\colon\Omega_XB_XA\to A\).
The ordinary statement is Loday--Vallette
\cite[Theorem~2.2.9]{LV12}; the displayed chiral statement is
conditional on the Ran comparison package used by Theorem~A.
\end{lemma}

\begin{proof}
The two bijections are represented by the same twisting morphism.
The identity of \(\Omega_XC\) represents \(\eta_C\), and the identity
of \(B_XA\) represents \(\varepsilon_A\).  Naturality gives the two
triangle identities.
\end{proof}

\begin{remark}[Three comparison maps;
\ClaimStatusDefinitional]
\label{rem:two-theorem-B-statements}
The adjunction supplies the universal maps
\[
 \eta_C\colon C\longrightarrow B_X\Omega_XC,
 \qquad
 \varepsilon_A\colon\Omega_XB_XA\longrightarrow A.
\]
A quadratic presentation supplies the additional map
\[
 q_A\colon A^{\mathrm i}\longrightarrow B_X(A).
\]
Theorem~A governs \(\eta_C\) and \(\varepsilon_A\) in the enhanced
Ran reconstruction ambient.  Theorem~B governs \(q_A\) and its cobar
mate \(\Omega_X(A^{\mathrm i})\to A\).  This assignment is stable
under every change of presentation and completion used below.
\end{remark}

\begin{theorem}[Universal coalgebra-side resolution;
\ClaimStatusConditional]
\label{thm:thm-B-coalgebra-side-unit-qi}
\textup{[Type signature: Open quadrant; conilpotent ordinary dg
coalgebra presentation; Beilinson levels \(2\leftrightarrow1\);
hypothesis package: ordinary conilpotence; pro-nilpotent chiral Ran
ambient with completed chain comparison \(H_{\mathrm{conv}}\).]}
Let \(C\) be a conilpotent coaugmented dg coalgebra.  The unit
\begin{equation}\label{eq:thm-B-coalg-unit}
 \eta_C\colon C\longrightarrow B\Omega(C)
\end{equation}
is a quasi-isomorphism
\cite[Corollary~2.3.4]{LV12}.  If \(C\) is a chiral coalgebra and
\(H_{\mathrm{conv}}\) identifies the ordinary filtered calculation
with the enhanced Ran bar and cobar constructions, the chiral unit
\(C\to B_X\Omega_XC\) is a quasi-isomorphism in that ambient.
\end{theorem}

\begin{proof}
The coradical filtration of \(C\) is exhaustive by conilpotence.
On associated graded objects, the unit is the standard tensor
coalgebra resolution.  Loday--Vallette's comparison lemma identifies
the unit as a quasi-isomorphism.  The chiral conclusion is transport
through the filtered comparison supplied by \(H_{\mathrm{conv}}\).
\end{proof}

\begin{remark}[Curved and second-kind ambient;
\ClaimStatusConditional]
\label{rem:thm-B-coalg-side-ambient}
A curved chiral CDG-coalgebra \(C_g=(C,d,h_g)\) carries the
fixed-coalgebra categories
\[
 D^{\mathrm{co}}(C_g\text{-}\mathrm{CoFact}),
 \qquad
 D^{\mathrm{ctr}}(C_g\text{-}\mathrm{ContraFact}).
\]
Their comparison is Positselski's fixed-\(C_g\) construction.  A coalgebra
map \(C_g\to B_X\Omega_XC_g\) in a curved model requires a curved
bar--cobar construction and a weak-equivalence structure on the
coalgebra category.  These data form an additional hypothesis
package.  The quadratic map \(q_A\) remains the carrier of
Theorem~B.
\end{remark}

\section{Theorem~B: quadratic Koszul recognition}
\label{sec:forward-references-thm-B-scope}

\begin{theorem}[Quadratic recognition in the chiral comparison
setting; \ClaimStatusConditional]
\label{thm:theorem-B-scope-quadratic-recognition}
\textup{[Type signature: Open quadrant; completed chiral
Ran/factorization presentation; Beilinson levels
\(1\leftrightarrow2\); hypothesis package: a connected positive
quadratic presentation and
\(H_{\mathrm{CL}}(A,A^{\mathrm i},\tau_{\mathrm i})\).]}
Let \(A=T_X(V)/(R)\), let
\(A^{\mathrm i}=C_X(s^{-1}V,s^{-2}R)\), and let
\(\tau_{\mathrm i}\colon A^{\mathrm i}\to A\) be the canonical
quadratic twisting morphism.  Under
\(H_{\mathrm{CL}}(A,A^{\mathrm i},\tau_{\mathrm i})\), the following
conditions are equivalent:
\begin{enumerate}[label=\textup{(\roman*)}]
\item \(\tau_{\mathrm i}\) is a Koszul twisting morphism;
\item \(q_A\colon A^{\mathrm i}\to B_X(A)\) is a quasi-isomorphism;
\item \(\Omega_X(A^{\mathrm i})\to A\) is a quasi-isomorphism;
\item the left and right twisted tensor products
 \(K^L_{\tau_{\mathrm i}}\) and \(K^R_{\tau_{\mathrm i}}\) are acyclic;
\item the completed PBW/bar spectral sequence is strongly convergent
 and concentrated on the quadratic Koszul diagonal.
\end{enumerate}
The ordinary equivalence of the twisting-morphism clauses is
Loday--Vallette~\cite[Theorem~2.3.2]{LV12}; its quadratic recognition
form is Loday--Vallette~\cite[Theorem~3.4.6]{LV12}.
\end{theorem}

\begin{proof}
Over a point, clauses \textup{(i)--(iv)} are the fundamental theorem
of twisting morphisms
\cite[Theorem~2.3.2]{LV12}, and clauses \textup{(ii)} and
\textup{(iii)} are the quadratic recognition theorem
\cite[Theorem~3.4.6]{LV12}.  The diagonal condition is the standard
filtered form of the same criterion.  The package
\(H_{\mathrm{CL}}\) transports these equivalences through chiral
strata, factorization gluing, and the strict completed limit.
\end{proof}

The completion results of
Theorem~\ref{thm:chiral-positselski-weight-completed} concern a fixed
coalgebra \(C\).  They can support the ambient construction of
\(B_X(A)\), while the proof of Theorem~B still requires the
quasi-isomorphism \(q_A\).  Likewise, a module-level Positselski
equivalence can transport modules attached to an acyclic twisting
datum, while the quadratic recognition problem remains the
acyclicity of the twisting datum
\((A,A^{\mathrm i},\tau_{\mathrm i})\).

\section{Arithmetic recognition over a base}
\label{sec:tbsp-supplement}

An arithmetic family contributes to Theorem~B through a family of
quadratic presentations and comparison maps.  Let \(S\) be a base and
let
\[
 A_S=T_X(V_S)/(R_S),\qquad
 A_S^{\mathrm i}=C_X(s^{-1}V_S,s^{-2}R_S),\qquad
 q_{A_S}\colon A_S^{\mathrm i}\longrightarrow B_X(A_S).
\]
A recognition datum over \(S\) consists of the presentation, the chiral
Comparison Lemma package, strong finite-window or completed
convergence, and a calculation of
\(\operatorname{Cone}(q_{A_S})\).  Automorphic, Hall, or K3 data may
supply this datum through an explicit filtered comparison
with \(q_{A_S}\).

\begin{remark}[K3 specialization as an external recognition datum;
\ClaimStatusOpen]
\label{rem:tbsp-1-ii}
A proposed K3 Hall/CoHA realization may be evaluated at a chosen
deformation parameter once source recognition, pairing,
PBW presentation completeness, parity, completion, inverse-limit, and
Heegner comparison data have been supplied.  The resulting object
enters Theorem~B through its displayed map \(q_A\).
\end{remark}

\begin{remark}[Arithmetic equations and the comparison cone;
\ClaimStatusOpen]
\label{rem:tbsp-2-ii}
An equation involving \(\hbar\), a conductor, or a modular divisor is
an arithmetic recognition condition.  Its connection with quadratic
Koszulness becomes a theorem after one constructs a morphism from the
arithmetic class to
\[
 H^\bullet\operatorname{Cone}
 \bigl(q_A\colon A^{\mathrm i}\to B_X(A)\bigr)
\]
and proves that the morphism detects vanishing.  This comparison is
an open obligation for the K3 surface discussed here.
\end{remark}

\section{The Etingof--Kazhdan completion as a comparison problem}
\label{sec:tbsp-etingof}

\begin{remark}[Two completed filtrations;
\ClaimStatusConditional]
\label{rem:tbsp-etingof-1}
\index{Etingof--Kazhdan!completion comparison}
The bar-weight completion of \(B_X(A)\) and the
\(\hbar\)-adic completion of an Etingof--Kazhdan quantization are two
filtered objects.  A filtered quasi-isomorphism between their Rees
objects supplies the bridge needed to compare them.  Theorem~B then
asks whether the quadratic subcoalgebra \(A^{\mathrm i}\) maps
quasi-isomorphically to the bar side.
\end{remark}

\begin{remark}[Order-\(\hbar^3\) target cocycle;
\ClaimStatusConditional]
\label{rem:tbsp-etingof-2}
\index{Siegel--Borcherds target cocycle!scope}
The candidate Siegel--Borcherds target cocycle records an
order-\(\hbar^3\) target-cocycle calculation in a completed
Chevalley--Eilenberg complex.  An all-order Hall associator
additionally requires a non-scalar source, an all-order pentagon
equation, and a filtered comparison with
\(B_X(A)\).  Once supplied, those data may be tested against
\(q_A\).
\end{remark}

\begin{remark}[Deformation-parameter convention;
\ClaimStatusOpen]
\label{rem:tbsp-etingof-3}
A Drinfeld deformation parameter, a root-of-unity parameter, and a
physical loop-counting parameter become one coordinate through an
explicit normalization theorem.  The K3 arithmetic surface
currently treats that normalization as recognition data.  Theorem~B
uses the resulting algebra and its quadratic comparison, rather than
a scalar equation by itself.
\end{remark}

\section{Quadratic recognition over a stratified base}
\label{sec:tbsp-supplement-ii}

Let \(S\) be a reduced base and let \(A_S=T(V_S)/(R_S)\) be a flat
family of complete augmented quadratic chiral algebras.  Put
\[
 K_{A_S}:=\operatorname{Cone}
 \bigl(q_{A_S}\colon A_S^{\mathrm i}\to B_X(A_S)\bigr).
\]
The quadratic Koszul locus is the open locus on which \(K_{A_S}\) is
acyclic whenever \(K_{A_S}\) is a perfect family.  Its complement is
stratified by the Fitting ideals of the cohomology sheaves of
\(K_{A_S}\).

\begin{remark}[Open quadratic Koszul locus;
\ClaimStatusConditional]
\label{rem:tbsp-1}
After a K3 Hall/CoHA candidate has acquired a quadratic presentation,
the open locus \(U\subset S\) belongs to Theorem~B precisely when
\(q_{A_S}|_U\) is a quasi-isomorphism under
\(H_{\mathrm{CL}}|_U\).  A divisor description of \(S\setminus U\)
requires an equality between those divisors and the Fitting support
of \(H^\bullet(K_{A_S})\).
\end{remark}

\begin{remark}[Completed quadratic comparison;
\ClaimStatusConditional]
\label{rem:tbsp-2}
Suppose \(A_S^{\mathrm i}\), \(B_X(A_S)\), and \(q_{A_S}\) are strict
inverse systems of finite windows and their cone tower is
Mittag--Leffler.  Windowwise quasi-isomorphisms then assemble to
\[
 \widehat q_{A_S}\colon
 \widehat{A_S^{\mathrm i}}\xrightarrow{\sim}\widehat{B_X(A_S)}.
\]
This is the completed form of Theorem~B.  The fixed-coalgebra
co--contra comparison of
Theorem~\ref{thm:chiral-positselski-weight-completed} remains a
separate categorical statement.
\end{remark}

\begin{remark}[Divisorial support of the comparison cone;
\ClaimStatusConditional]
\label{rem:tbsp-3}
A Humbert divisor \(H_n\) is a quadratic Koszul wall when
\[
 H_n\subset
 \operatorname{Supp}H^\bullet(K_{A_S}).
\]
Its multiplicity is the generic length of the relevant cohomology
sheaf along \(H_n\).  Monodromy, nearby cycles, and automorphic
coefficients can compute this length after a comparison theorem
identifies their filtered complexes with \(K_{A_S}\).
\end{remark}

\begin{corollary}[Local forms of the quadratic recognition problem;
\ClaimStatusConditional]
\label{cor:tbsp-trichotomy-HDelta5}
Assume a family \(A_S\) carries \(H_{\mathrm{CL}}\) and a perfect
comparison cone \(K_{A_S}\).  Then:
\begin{enumerate}[label=\textup{(\roman*)}]
\item on \(S\setminus\operatorname{Supp}H^\bullet(K_{A_S})\),
 \(q_{A_S}\) is a quasi-isomorphism;
\item on a formal divisor neighbourhood, Theorem~B is the vanishing
 problem for the completed cone;
\item on intersections of strata, global recognition is the descent
 problem for the local comparison maps.
\end{enumerate}
\end{corollary}

\begin{remark}[Higher operations as obstructions;
\ClaimStatusConditional]
\label{rem:tbsp-4}
A finite-window contraction of \(B_X(A)\) transfers its coalgebra
structure to a minimal homotopy-coalgebra structure on cohomology.
After a graded identification
\(A^{\mathrm i}\cong H^\bullet(B_X(A))\) is supplied, the first
nonzero higher cooperation obstructs strictification of this graded
identification to \(q_A\).  The cohomology of \(K_A\) records the
underlying quasi-isomorphism obstruction.  A numerical loop
coefficient enters this comparison through a chain-level
identification with the corresponding transferred cooperation.
\end{remark}

\begin{remark}[Local obstruction group;
\ClaimStatusConditional]
\label{rem:tbsp-5-tor-class}
On the formal completion \(\widehat S_{H_n}\), a finite flat
presentation of \(K_{A_S}\) defines local obstruction classes in its
cohomology and, after choosing a two-term resolution, in the
corresponding \(\operatorname{Tor}_1\) group.  An identification with
a Heegner, pentagon, or gerbe class is an additional comparison
theorem.  The present chapter records this as an open arithmetic
input.
\end{remark}

\begin{remark}[Per-window Mittag--Leffler criterion;
\ClaimStatusConditional]
\label{rem:tbsp-6-ml-per-weight}
For a strict tower \(K_{\leq N}\) of finite-window comparison cones,
the completed comparison is a quasi-isomorphism when
\[
 H^\bullet(K_{\leq N})=0
 \quad\text{for all }N,
 \qquad
 \varprojlim\nolimits^1_NH^{\bullet-1}(K_{\leq N})=0.
\]
Surjective transition maps give the standard sufficient
Mittag--Leffler condition.
\end{remark}

\begin{remark}[\(\mathcal D\)-module realization;
\ClaimStatusConditional]
\label{rem:tbsp-7-dmod-vs-positselski}
A regular holonomic realization of \(K_{A_S}\) carries the
Kashiwara--Malgrange \(V\)-filtration along a smooth divisor.  A
filtered Riemann--Hilbert comparison identifies its nearby cycles
with the completed algebraic cone.  This comparison can transport
monodromy calculations to Theorem~B.
\end{remark}

\begin{remark}[Admissible divisors as a proposed support;
\ClaimStatusOpen]
\label{rem:tbsp-8-admissible-extension}
A collection of admissible Humbert divisors defines a candidate
support for \(H^\bullet(K_{A_S})\).  Equality with the actual support
requires both
inclusions between this divisor union and the Fitting support of the
comparison cone.
\end{remark}

\begin{theorem}[Admissible-divisor criterion for quadratic
recognition; \ClaimStatusConditional]
\label{thm:tbsp-admissible-strict-scope}
\textup{[Type signature: Open quadrant over a stratified base; flat
quadratic chiral presentation; Beilinson levels
\(1\leftrightarrow2\); hypothesis package:
\(H_{\mathrm{CL}}\) and a perfect comparison cone.]}
Let \(S\) carry a flat quadratic family \(A_S\), its comparison
\(q_{A_S}\), and \(H_{\mathrm{CL}}\).  Let
\(D=\bigcup_{n\in I}H_n\) be a locally finite divisor union.  Assume:
\begin{enumerate}[label=\textup{(A\arabic*)}]
\item \(K_{A_S}\) is perfect;
\item \(H^\bullet(K_{A_S})\) is supported on \(D\);
\item the restriction of \(K_{A_S}\) to \(U=S\setminus D\) is
 acyclic.
\end{enumerate}
Then \(q_{A_S}|_U\) is a quasi-isomorphism.  Equality
\(D=\operatorname{Supp}H^\bullet(K_{A_S})\) follows when every
component \(H_n\) has a nonzero generic obstruction class.
\end{theorem}

\begin{proof}
Perfectness makes acyclicity an open condition.  Assumption
\textup{(A3)} gives the quasi-isomorphism on \(U\).  The final
assertion follows from the definition of scheme-theoretic support at
the generic point of each component.
\end{proof}

\section{Monodromy-refined finite-window comparison}
\label{sec:tbsp-strict-across-humbert-walls}

\begin{definition}[Monodromy-refined total bar weight;
\ClaimStatusDefinitional]
\label{def:monodromy-refined-total-bar-weight}
Let \(H\subset S\) be a smooth divisor, let \(T_H\) be the local
monodromy acting on the finite-window bar complex, and let
\(F^{\leq w}\) be total bar weight.  The refined filtration is
\[
 F^{\leq w,\chi}_H B_X(A)
 :=
 F^{\leq w}B_X(A)\cap B_X(A)^\chi .
\]
When the quadratic presentation is monodromy-equivariant, the same
filtration exists on \(A^{\mathrm i}\) and on
\(K_A=\operatorname{Cone}(q_A)\).
\end{definition}

\begin{lemma}[Equivariance of the finite-window bar differential;
\ClaimStatusConditional]
\label{lem:bar-differential-preserves-monodromy-filtration}
Assume the local system, quadratic presentation, and chiral
operations are monodromy-equivariant.  Then the bar differential and
\(q_A\) preserve \(F^{\leq w,\chi}_H\).
\end{lemma}

\begin{proof}
On each genus-\(0\) affine/tangent finite KZ window, the singular
collision term is the Fulton--MacPherson boundary-residue realisation
of the KZ--Arnold bar superconnection.  This finite-window statement
relates objects of distinct types through the Fulton--MacPherson
residue-realization functor.  Horizontality gives the
finite-window flat-connection statement.  The internal and residue
terms are equivariant by the stated hypothesis, and the quadratic
relations are stable under monodromy.  Thus both differentials and
\(q_A\) preserve the refined filtration.
\end{proof}

\begin{lemma}[Per-isotypic Mittag--Leffler criterion;
\ClaimStatusConditional]
\label{lem:per-isotypic-ML-Hn}
Suppose each finite-window isotypic cone
\(K_{\leq w,\chi}\) has finite-dimensional cohomology and the images
of
\[
 H^m(K_{\leq w+1,\chi})\longrightarrow
 H^m(K_{\leq w,\chi})
\]
stabilize for every \((m,\chi)\).  Then
\(\varprojlim^1_wH^{m-1}(K_{\leq w,\chi})=0\), and the completed
cohomology is the inverse limit of the finite-window cohomologies.
\end{lemma}

\begin{theorem}[Monodromy-refined quadratic comparison;
\ClaimStatusConditional]
\label{thm:strict-chain-level-across-humbert-walls}
\textup{[Type signature: Open quadrant on a formal divisor
neighbourhood; monodromy-refined completed quadratic presentation;
Beilinson levels \(1\leftrightarrow2\); hypothesis package:
\(H_{\mathrm{CL}}\), windowwise recognition, and per-isotypic
Mittag--Leffler.]}
Let \(A_{\widehat S_H}\) be a monodromy-equivariant quadratic family
on a formal divisor neighbourhood.  Assume
\(H_{\mathrm{CL}}\), finite-window quasi-isomorphisms
\(q_{\leq w,\chi}\), and the Mittag--Leffler condition of
Lemma~\ref{lem:per-isotypic-ML-Hn}.  Then
\[
 \widehat q_A^\Gamma\colon
 \widehat{A^{\mathrm i}}^\Gamma
 \xrightarrow{\sim}
 \widehat{B_X(A)}^\Gamma
\]
is a quasi-isomorphism in the monodromy-refined completed ambient.
If a finite-window strong deformation retract is supplied, its
homotopy satisfies
\begin{equation}\label{eq:chain-homotopy-identity-humbert}
 [d,h]=\mathrm{id}-\iota p .
\end{equation}
\end{theorem}

\begin{proof}
The cone of \(\widehat q_A^\Gamma\) is the derived inverse limit of
the finite-window cones.  Their cohomology vanishes, and the Milnor
\(\varprojlim^1\) term vanishes by the stated condition.  Hence the
completed cone is acyclic.

For a supplied contraction of the unperturbed finite-window complex,
write
\begin{equation}\label{eq:bar-shift-koszul-splitting}
 d_0h_0+h_0d_0=\mathrm{id}-\iota p .
\end{equation}
If a perturbation \(\delta\) is filtration-lowering and locally
nilpotent, the homological perturbation lemma gives
\begin{equation}\label{eq:curvature-corrected-bott-koszul}
 h_\delta=h_0(1+\delta h_0)^{-1}
 =\sum_{r\geq0}(-1)^rh_0(\delta h_0)^r .
\end{equation}
The square-zero condition for the perturbed differential is the
Maurer--Cartan equation
\begin{equation}\label{eq:mc-integrability-curvature-corrected}
 [d_0,\delta]+\tfrac12[\delta,\delta]=0 .
\end{equation}
These supplied identities yield the displayed strong deformation
retract.
\end{proof}

\begin{remark}[Local and global comparison data;
\ClaimStatusConditional]
\label{rem:monodromy-refined-trichotomy}
The local-to-global argument has three components: the open acyclicity of
\(K_A\), the completed monodromy-refined acyclicity on divisor
neighbourhoods, and descent on their intersections.  Each component is a
statement about the same map \(q_A\).
\end{remark}

\begin{theorem}[Wall-of-walls descent obstruction;
\ClaimStatusOpen]
\label{thm:wall-of-walls-obstruction}
\textup{[Type signature: Open quadrant on a codimension-two
intersection; derived comparison-cone presentation; Beilinson levels
\(1\leftrightarrow2\); hypothesis package: local quadratic
comparisons and their transition homotopies.]}
On an intersection \(H_i\cap H_j\), the obstruction to gluing two
local quadratic comparisons is the \v Cech class of their transition
automorphism in the derived endomorphism complex of
\(K_A|_{H_i\cap H_j}\).  An explicit formula in terms of monodromy
commutators, cup products, or torsion orders requires a comparison
from those geometric classes to this derived endomorphism complex.
Construction of this comparison remains an open obligation in this
chapter.
\end{theorem}

\begin{remark}[Derived descent across intersections;
\ClaimStatusConditional]
\label{rem:wall-of-walls-rescue}
A coherent null-homotopy of the \v Cech class gives descent of the
 local maps \(q_A\).  A coderived or pro-completed category contains
 the null-homotopy when it exists there.  The category and
the null-homotopy are part of the hypothesis package.
\end{remark}

\begin{remark}[Three tests of the local comparison;
\ClaimStatusConditional]
\label{rem:three-checks-strict-humbert}
A monodromy-refined recognition claim should be checked by: a direct
finite-window cone calculation; a filtered nearby-cycle comparison;
and a limiting calculation on the open quadratic Koszul locus.
Agreement of these three paths proves the local statement.
\end{remark}

\begin{corollary}[Stratified comparison package;
\ClaimStatusConditional]
\label{cor:sharpened-quadrichotomy-after-monodromy-refinement}
Under the hypotheses of
Theorems~\ref{thm:tbsp-admissible-strict-scope}
and~\ref{thm:strict-chain-level-across-humbert-walls}, together with a
coherent descent null-homotopy on every intersection, the family map
\(q_{A_S}\) is a quasi-isomorphism in the assembled completed ambient.
\end{corollary}

\begin{remark}[Higher monodromy refinements;
\ClaimStatusOpen]
\label{rem:pentachotomy-to-hexachotomy-bi-unipotent}
A Deligne--Malcev or multi-\(V\)-filtration refinement can supply the
descent null-homotopies at normal-crossing intersections.  Applying
such a refinement to the K3 candidate requires a theorem identifying
its filtered object with the comparison cone \(K_A\).
\end{remark}

\section{Nearby-cycle tests along admissible divisors}
\label{sec:tbsp-chain-level-nearby-cycle-admissible-Hn}

\subsection{The nearby-cycle \(\mathcal D\)-module}
\label{subsec:tbsp-nearby-cycle-D-module}

\begin{definition}[Nearby-cycle comparison object;
\ClaimStatusDefinitional]
\label{def:tbsp-chain-nearby-cycle}
For a regular holonomic realization of \(K_A\) along a smooth divisor
\(H_n\), define
\[
 \psi_n(K_A)
 :=
 \bigoplus_{\lambda\in[0,1)\cap\mathbb Q}
 \operatorname{gr}^V_\lambda(K_A)|_{H_n}.
\]
It carries the usual quasi-unipotent monodromy.
\end{definition}

\begin{lemma}[Unipotent contribution to the local cone;
\ClaimStatusConditional]
\label{lem:tbsp-unipotent-part-controls}
If every nontrivial monodromy eigensummand admits the standard
projector contraction, then the cohomology of the local completed
cone is supported on the generalized eigenvalue-\(1\) summand of
\(\psi_n(K_A)\).
\end{lemma}

\begin{remark}[Arithmetic admissibility;
\ClaimStatusOpen]
\label{rem:tbsp-why-admissibility}
A congruence condition on Humbert discriminants selects a geometric
divisor family.  Its equality with the support of
\(H^\bullet(K_A)\) is the arithmetic-to-Koszul comparison obligation.
\end{remark}

\subsection{The filtered quadratic cone}
\label{subsec:tbsp-koszul-bar-filtration}

\begin{definition}[\(V\)-filtration on the comparison cone;
\ClaimStatusDefinitional]
\label{def:tbsp-koszul-bar-filtration}
Set
\[
 F^kK_A:=V^{\geq k/\ell_n}_{H_n}K_A,
 \qquad
 \operatorname{gr}^k_FK_A
 =\operatorname{gr}^{V}_{k/\ell_n}K_A .
\]
This filtration compares the quadratic and bar sides simultaneously.
\end{definition}

\begin{lemma}[Finite-graded recognition;
\ClaimStatusConditional]
\label{lem:tbsp-positselski-on-each-Fk}
If
\(\operatorname{gr}^k_Fq_A\) is a quasi-isomorphism for every \(k\)
and the filtered comparison is strongly convergent, then \(q_A\) is a
quasi-isomorphism.  This is the filtered comparison lemma applied to
the cone \(K_A\).
\end{lemma}

\subsection{The Milnor obstruction}
\label{subsec:tbsp-lim1-obstruction}

\begin{theorem}[Arithmetic coefficient to Milnor obstruction;
\ClaimStatusOpen]
\label{thm:tbsp-lim1-is-fourier-coefficient}
\textup{[Type signature: Open quadrant along an arithmetic divisor;
filtered comparison-cone presentation; Beilinson levels
\(1\leftrightarrow2\); hypothesis package: a filtered
automorphic-to-bar comparison.]}
Let \(c(n,r_n)\) be an automorphic coefficient attached to \(H_n\).
An equality
\[
 [\operatorname{obs}^{\lim^1}_{H_n}]
 =
 c(n,r_n)[H_n]^\vee
 \quad\text{in}\quad
 \varprojlim\nolimits^1_k
 H^{\bullet-1}(\operatorname{gr}^k_FK_A)
\]
requires a filtered comparison identifying the Chern/residue class of
the automorphic object with the extension class of the inverse system
\(F^\bullet K_A\).  The equality becomes a theorem after that
comparison is constructed and checked on a normalization of \(H_n\).
\end{theorem}

\begin{remark}[Leading-wall calculation;
\ClaimStatusOpen]
\label{rem:tbsp-leading-three-walls-first-coefficients}
The proposed leading Humbert coefficients provide concrete inputs for
Theorem~\ref{thm:tbsp-lim1-is-fourier-coefficient}.  Their numerical
values, normalizations, and images in the Milnor group require direct
calculation from the chosen Jacobi form and the chain-level filtered
comparison.
\end{remark}

\subsection{Boundary compatibility}
\label{subsec:tbsp-baily-borel-boundary}

\begin{proposition}[Nearby cycles and Baily--Borel base change;
\ClaimStatusConditional]
\label{prop:tbsp-baily-borel-compatibility}
Assume the comparison cone extends as a mixed Hodge module to the
Baily--Borel compactification and is noncharacteristic for the
boundary strata.  Then Saito's nearby-cycle base-change theorem
identifies the restriction of \(\psi_n(K_A)\) with the nearby cycles
of the restricted comparison cone.  The Milnor obstruction is
functorial under this restriction \cite{Saito1988}.
\end{proposition}

\subsection{Global quadratic recognition}
\label{subsec:tbsp-global-inversion-theorem}

\begin{theorem}[Global quadratic Koszul recognition on a stratified
base; \ClaimStatusConditional]
\label{thm:tbsp-global-inversion-all-admissible}
\textup{[Type signature: Open quadrant over the genus-two modular
base \(\overline{\mathcal A_2}\); chiral bar chain presentation for
the target \(B_X(A)\), with a quadratic source \(A^{\mathrm i}\);
Beilinson level \(1\leftrightarrow2\); global assembly in the
completed chiral coalgebra ambient serving the descent step.  The hypotheses
named below are \(H_{\mathrm{CL}}\), perfect finite-window comparison
cones, strict Mittag--Leffler convergence, regular nearby-cycle
realization, and coherent \v Cech descent.]}

Let \(S=\overline{\mathcal A_2}\), let
\(A_S=T_X(V_S)/(R_S)\), and let
\(q_{A_S}\colon A_S^{\mathrm i}\to B_X(A_S)\).  Choose an open set
\(U\) and formal divisor neighbourhoods
\(\widehat S_{H_n}\) covering \(S\).  Assume:
\begin{enumerate}[label=\textup{(G\arabic*)}]
\item \(q_{A_S}|_U\) is a quasi-isomorphism;
\item each monodromy-refined finite-window map on
 \(\widehat S_{H_n}\) is a quasi-isomorphism;
\item the cone towers are strict Mittag--Leffler;
\item the local maps and their homotopies satisfy coherent \v Cech
 descent on all intersections.
\end{enumerate}
Then the assembled map
\[
 \widehat q_{A_S}\colon
 \widehat{A_S^{\mathrm i}}
 \xrightarrow{\sim}
 \widehat{B_X(A_S)}
\]
is a quasi-isomorphism.  Equivalently, \(A_S\) satisfies Theorem~B in
the stated completed chiral ambient.
\end{theorem}

\begin{proof}
Work on the \v Cech cover in the finite-window chiral chain model.
The target differential is the chiral bar differential on
\(B_X(A_S)\), the source carries the quadratic differential, and the
map is \(q_{A_S}\).  Filter every local comparison cone by the
monodromy-refined total bar-weight filtration.  The finite-window bar
spectral sequence and the quadratic source spectral sequence are
identified by the local hypotheses.  The strict Mittag--Leffler
hypothesis kills the Milnor \(\varprojlim^1\) term.  The coherent
descent data glue the acyclic local cones to an acyclic global cone.
\end{proof}

\begin{remark}[Tests of global recognition;
\ClaimStatusConditional]
\label{rem:tbsp-global-inversion-independent-checks}
A global recognition argument combines a direct calculation of
\(\operatorname{Cone}(q_A)\), a nearby-cycle calculation of its
filtered realization, and an open-locus PBW calculation.  A
cross-volume arithmetic invariant becomes a fourth check after its
comparison map to the cone has been constructed.
\end{remark}

\begin{remark}[Ambient-neutral form of Theorem~B;
\ClaimStatusConditional]
\label{rem:tbsp-ambient-neutral-form}
The ambient-neutral form of Theorem~B is the local-to-global statement
for \(q_A\): ordinary finite windows on the open locus, strict
completed windows near the boundary, and coherent descent on
intersections.  Each ambient realizes the same quadratic comparison
map.
\end{remark}

\section{Fixed-coalgebra and module comparisons}
\label{sec:tbsp-module-level-proof-valpha1}

The module layer adjacent to Theorem~B has two rigorous forms.  For a
fixed chiral CDG-coalgebra \(C\), Positselski's theorem compares
\(C\)-comodules and \(C\)-contramodules.  For a twisting morphism
\(\tau\colon C\to A\) satisfying the acyclicity package, the twisted
tensor functors compare \(C\)-comodules with \(A\)-modules.  These are
module-category theorems.  The quadratic specialization
\((C,\tau)=(A^{\mathrm i},\tau_{\mathrm i})\) is connected with
Theorem~B exactly when \(q_A\) is a quasi-isomorphism.

The normalized two-sided bar resolution supplies a third, elementary
module construction.  It resolves an \(A\)-module for every augmented
algebra \(A\); its chiral transport is conditional on the
finite-window Ran comparison.  We develop this model for a proposed
module \(V_{\alpha_1}\) while keeping its recognition hypotheses
visible.

\subsection{The proposed module \(V_{\alpha_1}\)}
\label{subsec:tbsp-module-valpha1}

\begin{definition}[Fundamental-root module datum;
\ClaimStatusDefinitional]
\label{def:tbsp-valpha1}
Assume a recognized Hall/BKM algebra \(A\) contains an
\(\mathfrak{sl}_2\)-triple
\((e_{\alpha_1},f_{\alpha_1},h_{\alpha_1})\).  The fundamental-root
module datum is
\begin{equation}\label{eq:tbsp-valpha1-definition}
 V_{\alpha_1}=\mathbb C v^+\oplus\mathbb C v^-,
 \qquad
 e_{\alpha_1}v^-=v^+,\quad
 f_{\alpha_1}v^+=v^-,\quad
 h_{\alpha_1}v^\pm=\pm v^\pm .
\end{equation}
An \(A\)-module or \(A\)-comodule structure extending this
\(\mathfrak{sl}_2\)-action is additional recognition data.
\end{definition}

\begin{remark}[Recognition obligation for the chiral action;
\ClaimStatusOpen]
\label{rem:tbsp-valpha1-legitimate}
A Hall/CoHA recognition package consists of the action map
\(A\otimes V_{\alpha_1}\to V_{\alpha_1}\), its chiral locality, and
its compatibility with the chosen completion.  A comodule
realization instead requires a coaction and the corresponding fixed
coalgebra.  The formulas below use the module realization.
\end{remark}

\subsection{The module bar construction}
\label{subsec:tbsp-barvalpha1}

\begin{definition}[Normalized one-sided module bar complex;
\ClaimStatusDefinitional]
\label{def:tbsp-barch-module}
For an augmented chiral algebra \(A\) and a right \(A\)-module \(M\),
define
\begin{equation}\label{eq:tbsp-barch-module}
 B_X^{\mathrm{ch}}(M;A;\mathbf1)
 :=
 M\otimes T_X^c(s^{-1}\bar A)
 =
 \bigoplus_{p\geq0}M\otimes(s^{-1}\bar A)^{\otimes p}.
\end{equation}
The differential is the sum of the internal differential, adjacent
multiplications in \(A\), and the module-action face at the \(M\)-end.
Suspension conventions are those of
Loday--Vallette~\cite[\S2.2]{LV12}.
\end{definition}

\begin{lemma}[First three bar degrees;
\ClaimStatusProvedHere]
\label{lem:tbsp-bar-valpha1-first-terms}
For \(M=V_{\alpha_1}\),
\begin{align}
 B_X^{\mathrm{ch},0}(V_{\alpha_1};A;\mathbf1)
 &=V_{\alpha_1},
 \label{eq:tbsp-bar-valpha1-deg0}\\
 B_X^{\mathrm{ch},1}(V_{\alpha_1};A;\mathbf1)
 &=V_{\alpha_1}\otimes s^{-1}\bar A,
 \label{eq:tbsp-bar-valpha1-deg1}\\
 B_X^{\mathrm{ch},2}(V_{\alpha_1};A;\mathbf1)
 &=V_{\alpha_1}\otimes(s^{-1}\bar A)^{\otimes2}.
 \label{eq:tbsp-bar-valpha1-deg2}
\end{align}
If \(A\) has finite-dimensional positive-energy weight spaces, every
fixed total-weight piece of these three terms is finite-dimensional.
\end{lemma}

\begin{proof}
The formulas are the components \(p=0,1,2\) of
Definition~\ref{def:tbsp-barch-module}.  Finite-dimensionality follows
from the finite number of weight decompositions in a fixed total
weight.
\end{proof}

\subsection{The two-sided resolution}
\label{subsec:tbsp-omegabar-valpha1}

\begin{definition}[Normalized two-sided module model;
\ClaimStatusDefinitional]
\label{def:tbsp-omegabar-module}
For a right \(A\)-module \(M\), set
\begin{equation}\label{eq:tbsp-two-sided-model}
 K_X^{\mathrm{ch}}(A,M)
 :=
 M\otimes T_X^c(s^{-1}\bar A)\otimes A .
\end{equation}
The differential contains the internal differentials, adjacent
products, the module action, and multiplication into the terminal
\(A\)-factor.  A chiral comparison map
\begin{equation}\label{eq:tbsp-two-sided-comparison}
 \theta_M\colon
 K_X^{\mathrm{ch}}(A,M)
 \longrightarrow
 \Omega_X^{\mathrm{ch}}B_X^{\mathrm{ch}}(M;A;\mathbf1)
\end{equation}
requires a typed chiral twisting construction.  Its existence and
filtered quasi-isomorphism property form the comparison package
\(H_{\mathrm{mod}}(A,M)\).
\end{definition}

\begin{remark}[The terminal algebra factor]
\label{rem:tbsp-two-sided-vs-bicomplex}
The terminal \(A\)-factor carries the standard extra degeneracy.
The distinguished module vector remains in its \(M\)-slot throughout
the contraction.  This typing is the structural reason for using the
two-sided model.
\end{remark}

\begin{remark}[Finite-window totalization;
\ClaimStatusConditional]
\label{rem:tbsp-coderived-total}
In a fixed conformal-weight window the two-sided model is finite in
each total degree.  Product or inverse-limit totalization enters when
a completed tower is assembled.  The strict Mittag--Leffler
package controls that passage.
\end{remark}

\subsection{Projection, inclusion, and contraction}
\label{subsec:tbsp-homotopy-valpha1}

\begin{construction}[Two-sided contraction]
\label{cons:tbsp-projection-homotopy}
Let \(a=\bar a+\epsilon(a)1\).  The augmentation projection and
inclusion are
\[
 \pi_K(m\otimes[]\otimes a)=m\cdot a,
 \qquad
 \iota_K(m)=m\otimes[]\otimes1.
\]
The standard extra degeneracy is
\begin{equation}\label{eq:tbsp-homotopy-formula}
 H_K\bigl(
 m\otimes[s^{-1}\bar a_1|\cdots|s^{-1}\bar a_p]\otimes a
 \bigr)
 =
 (-1)^\sigma
 m\otimes[s^{-1}\bar a_1|\cdots|s^{-1}\bar a_p|
 s^{-1}\bar a]\otimes1,
\end{equation}
with the Koszul sign determined by the suspension convention.
\end{construction}

\begin{theorem}[Two-sided reconstruction for
\(V_{\alpha_1}\); \ClaimStatusConditional]
\label{thm:tbsp-module-theorem-B-valpha1}
\textup{[Type signature: Open quadrant; finite-window right-module
two-sided bar presentation; Beilinson level~\(1\) module category;
hypothesis package: a chiral \(A\)-action and
\(H_{\mathrm{mod}}(A,V_{\alpha_1})\).]}
Assume \(V_{\alpha_1}\) is a right \(A\)-module in the finite-window
chiral model.  Then
\[
 \pi_K\colon K_X^{\mathrm{ch}}(A,V_{\alpha_1})
 \longrightarrow V_{\alpha_1}
\]
is a chain-homotopy equivalence, with
\begin{equation}\label{eq:tbsp-homotopy-identity}
 D_KH_K+H_KD_K
 =
 \mathrm{id}-\iota_K\pi_K .
\end{equation}
Under \(H_{\mathrm{mod}}(A,V_{\alpha_1})\), the transported counit
\begin{equation}\label{eq:tbsp-module-counit-chain}
 \Omega_X^{\mathrm{ch}}
 B_X^{\mathrm{ch}}(V_{\alpha_1};A;\mathbf1)
 \xrightarrow{\sim}V_{\alpha_1}
\end{equation}
is a quasi-isomorphism.  This is a module reconstruction statement;
Theorem~B enters after the twisting coalgebra is specialized to
\(A^{\mathrm i}\) and \(q_A\) is proved to be a quasi-isomorphism.
\end{theorem}

\begin{proof}
The normalized two-sided bar complex has the extra degeneracy
\eqref{eq:tbsp-homotopy-formula}.  The adjacent-product terms cancel
pairwise with the terminal multiplication and module-action terms,
leaving the module augmentation \(\iota_K\pi_K\).  Transport through
\(\theta_{V_{\alpha_1}}\) follows from two-out-of-three under
\(H_{\mathrm{mod}}\).
\end{proof}

\begin{corollary}[Named-module reconstruction;
\ClaimStatusConditional]
\label{cor:tbsp-named-module-theorem-B}
The explicit two-sided model proves module reconstruction for
\(V_{\alpha_1}\) once its \(A\)-action is supplied.  Transport to a
chiral cobar-of-bar presentation requires \(H_{\mathrm{mod}}\).
Transport through a quadratic twisting datum additionally requires
Theorem~B for \(q_A\).
\end{corollary}

\begin{remark}[Three tests of the module statement]
\label{rem:tbsp-module-four-checks-valpha1}
The module statement has three independent checks: the extra
degeneracy on the two-sided model, collapse of its bar spectral
sequence, and specialization to an ordinary augmented algebra over a
point.  A family-specific Hall or BKM action supplies a fourth
recognition check.
\end{remark}

\begin{remark}[Extension to a Mukai module;
\ClaimStatusOpen]
\label{rem:tbsp-heisenberg-yangian-standard}
A proposed Mukai-lattice module requires an action of the recognized
algebra on \(H^\bullet(K3)\), compatibility with the \(24\)-dimensional
Mukai grading, and a continuous finite-window decomposition.  The
two-sided reconstruction then applies to the full module directly;
a decomposition into root directions is optional auxiliary data.
\end{remark}

\begin{remark}[Divisorwise module transport;
\ClaimStatusConditional]
\label{rem:tbsp-module-humbert-scope}
If \(A_S\), \(M_S\), and \(H_{\mathrm{mod}}(A_S,M_S)\) vary over a
stratified base, the two-sided contraction is natural under base
change.  Completed descent requires strict Mittag--Leffler control of
the module cone tower.  A quadratic wall affects the module category
through the twisting comparison \(q_{A_S}\).
\end{remark}

\section{All bar degrees and the two-sided spectral sequence}
\label{sec:tbsp-higher-bar-degrees-valpha1}

\subsection{Bar spectral sequence}
\label{subsec:tbsp-bar-ss-valpha1}

\begin{definition}[Bar filtration on the two-sided model;
\ClaimStatusDefinitional]
\label{def:tbsp-bar-spectral-sequence}
Filter \(K_X^{\mathrm{ch}}(A,M)\) by bar length:
\[
 F_pK=\bigoplus_{0\leq r\leq p}
 M\otimes(s^{-1}\bar A)^{\otimes r}\otimes A .
\]
In each finite weight window this filtration is bounded and
exhaustive.
\end{definition}

\begin{lemma}[Collapse of the two-sided bar spectral sequence;
\ClaimStatusProvedElsewhere]
\label{lem:tbsp-bar-ss-collapse}
The spectral sequence has
\[
 E_2^{p,q}\cong
 \begin{cases}
 H^q(M),&p=0,\\
 0,&p>0,
 \end{cases}
\]
and converges to \(H^\bullet(M)\).  Under \(H_{\mathrm{mod}}(A,M)\)
the same conclusion transports to the displayed chiral
cobar-of-bar model.
\end{lemma}

\begin{proof}
The extra degeneracy contracts every positive bar-length column.
The augmentation column is \(M\).  Finite-window boundedness gives
strong convergence.
\end{proof}

\begin{remark}[Structural control of Koszul signs]
\label{rem:tbsp-exponential-signs}
The extra-degeneracy identity contains the complete Koszul sign
calculation in every bar degree.  The spectral sequence records the
same contraction at the level of filtered homology.
\end{remark}

\subsection{Selected bar degrees}
\label{subsec:tbsp-n456-valpha1}

\begin{proposition}[Bar degree \(4\);
\ClaimStatusProvedElsewhere]
\label{prop:tbsp-homotopy-n4-valpha1}
On
\(M\otimes(s^{-1}\bar A)^{\otimes4}\otimes A\), the two-sided homotopy
satisfies
\begin{equation}\label{eq:tbsp-homotopy-n4}
 D_KH_K+H_KD_K=\mathrm{id}-\iota_K\pi_K .
\end{equation}
\end{proposition}

\begin{proof}
This is the degree-\(4\) component of the normalized two-sided bar
contraction.
\end{proof}

\begin{proposition}[Bar degree \(5\) under a filtered perturbation;
\ClaimStatusConditional]
\label{prop:tbsp-homotopy-n5-valpha1}
Let \(\delta\) be a locally nilpotent filtered perturbation compatible
with the \(A\)-module structure.  The homological perturbation lemma
produces \(H_\delta\) and
\begin{equation}\label{eq:tbsp-homotopy-n5}
 (D_K+\delta)H_\delta+H_\delta(D_K+\delta)
 =\mathrm{id}-\iota_\delta\pi_\delta
\end{equation}
in bar degree \(5\).
\end{proposition}

\begin{proof}
Apply the finite geometric-series formula of
\eqref{eq:curvature-corrected-bott-koszul} in the fixed weight window.
\end{proof}

\begin{proposition}[Bar degree \(6\);
\ClaimStatusProvedElsewhere]
\label{prop:tbsp-homotopy-n6-valpha1}
The unperturbed two-sided contraction satisfies the same identity in
bar degree \(6\).
\end{proposition}

\begin{proof}
This is the degree-\(6\) component of the extra-degeneracy identity.
\end{proof}

\begin{corollary}[Uniform two-sided contraction;
\ClaimStatusProvedElsewhere]
\label{cor:tbsp-homotopy-all-n-valpha1}
The identity
\(D_KH_K+H_KD_K=\mathrm{id}-\iota_K\pi_K\)
holds in every bar degree.
\end{corollary}

\section{Additional module test objects}
\label{sec:tbsp-additional-comodules}

\subsection{Vacuum module}
\label{subsec:tbsp-vacuum-module}

\begin{definition}[Vacuum module]
\label{def:tbsp-vacuum-comodule}
The vacuum is the one-dimensional module \(\mathbf1\) obtained from
the augmentation \(A\to\mathbb C\).
\end{definition}

\begin{theorem}[Two-sided reconstruction of the vacuum;
\ClaimStatusConditional]
\label{thm:tbsp-theorem-B-vacuum}
\textup{[Type signature: Open quadrant; finite-window augmentation
module presentation; Beilinson level~\(1\) module category; hypothesis
package: the chiral two-sided bar comparison.]}
Under the finite-window chiral comparison,
\begin{equation}\label{eq:tbsp-theorem-B-vacuum}
 K_X^{\mathrm{ch}}(A,\mathbf1)
 \xrightarrow{\sim}\mathbf1
\end{equation}
is a quasi-isomorphism.
\end{theorem}

\begin{proof}
Apply Theorem~\ref{thm:tbsp-module-theorem-B-valpha1} with
\(M=\mathbf1\).
\end{proof}

\begin{remark}[Vacuum as an augmentation check]
\label{rem:tbsp-vacuum-base-case}
The vacuum tests the augmentation and the terminal multiplication in
the two-sided model.  It is a module-level check on Theorem~A's
reconstruction map.
\end{remark}

\subsection{Adjoint module}
\label{subsec:tbsp-adjoint-module}

\begin{definition}[Completed adjoint-module datum;
\ClaimStatusDefinitional]
\label{def:tbsp-adjoint-module}
For a recognized topological Lie or BKM algebra \(\mathfrak g\), let
\(\widehat{\operatorname{ad}}(\mathfrak g)\) denote a supplied
weight-completed adjoint module whose finite windows are
finite-dimensional and whose transition maps are continuous.
\end{definition}

\begin{theorem}[Two-sided reconstruction of a completed adjoint
module; \ClaimStatusConditional]
\label{thm:tbsp-theorem-B-adjoint}
\textup{[Type signature: Open quadrant; strict completed adjoint-module
presentation; Beilinson level~\(1\) module category; hypothesis package:
continuous action, finite-window comparison, and Mittag--Leffler.]}
Assume the adjoint action extends to a continuous \(A\)-module
structure and that its finite-window towers satisfy strict
Mittag--Leffler.  Then the windowwise maps
\begin{equation}\label{eq:tbsp-theorem-B-adjoint}
 K_X^{\mathrm{ch}}(A,\operatorname{ad}_{\leq N})
 \xrightarrow{\sim}\operatorname{ad}_{\leq N}
\end{equation}
assemble to a quasi-isomorphism on the completed adjoint module.
\end{theorem}

\begin{proof}
Use the two-sided contraction at every finite window and the Milnor
exact sequence for the strict tower.
\end{proof}

\begin{remark}[Completion of the adjoint module]
\label{rem:tbsp-adjoint-weight-completion}
The completed statement is governed by continuity of the action and
vanishing of the \(\varprojlim^1\) term.  Root multiplicity formulas
establish finite-window dimensions and transition maps.
\end{remark}

\subsection{An isotropic-root test module}
\label{subsec:tbsp-isotropic-imaginary-valpha}

\begin{definition}[Isotropic-root module datum;
\ClaimStatusDefinitional]
\label{def:tbsp-vbeta-isotropic}
Let \(V_\beta\) be a supplied finite-dimensional module attached to
an isotropic root \(\beta\), together with a continuous \(A\)-action.
For a \(20\)-dimensional slice, its geometric source is specified
explicitly: \(h^{1,1}(K3)=20\), whereas
\(\chi_{\mathrm{top}}(K3)=24\).
\end{definition}

\begin{theorem}[Two-sided reconstruction of \(V_\beta\);
\ClaimStatusConditional]
\label{thm:tbsp-theorem-B-Vbeta}
\textup{[Type signature: Open quadrant; finite-window isotropic-root
module presentation; Beilinson level~\(1\) module category; hypothesis
package: a continuous action and the chiral two-sided bar comparison.]}
Under the finite-window chiral module comparison,
\begin{equation}\label{eq:tbsp-theorem-B-Vbeta}
 K_X^{\mathrm{ch}}(A,V_\beta)
 \xrightarrow{\sim}V_\beta
\end{equation}
is a quasi-isomorphism.  A completed family statement follows from
strict Mittag--Leffler control.
\end{theorem}

\begin{proof}
Apply the normalized two-sided contraction, followed by the finite
window or strict inverse-limit comparison.
\end{proof}

\begin{remark}[Module objects adjacent to Theorem~B]
\label{rem:tbsp-named-comodules-module-surface}
The vacuum, fundamental-root, adjoint, and isotropic-root modules test
module reconstruction and the fixed-coalgebra comparison.  They become
tests of Theorem~B when the relevant twisting coalgebra is
\(A^{\mathrm i}\) and the comparison
\(q_A\colon A^{\mathrm i}\to B_X(A)\) is a quasi-isomorphism.
\end{remark}
